\chapter{Cat\'egories fibr\'ees et descente}
\label{VI}

\setcounter{section}{-1}
\section{Introduction}

Contrairement \`a ce qui avait \'et\'e annonc\'e dans
l'introduction \`a l'expos\'e pr\'ec\'edent, il s'est
av\'er\'e impossible de faire de la descente dans la cat\'egorie
des pr\'esch\'emas, m\^eme dans des cas particuliers, sans avoir
d\'evelopp\'e au pr\'ealable avec assez de soin le langage de la
descente dans les cat\'egories g\'en\'erales.

La notion de \og descente\fg fournit le cadre g\'en\'eral pour tous
les proc\'ed\'es de \og recollement\fg d'objets, et par
cons\'equent de \og recollement\fg de cat\'egories. Le cas le plus
classique de recollement est relatif \`a la donn\'ee d'un espace
topologique $X$ et d'un recouvrement de~$X$ par des ouverts $X_i$; si
on se donne pour tout $i$ un espace fibr\'e (disons) $E_i$ au-dessus
de~$X_i$, et pour tout couple $(i,j)$ un isomorphisme $f_{ji}$ de
$E_i|X_{ij}$ sur~$E_j|X_{ij}$ (o\`u on pose $X_{ij}=X_i\cap X_j$),
satisfaisant une condition de transitivit\'e bien connue (qu'on
\'ecrit de fa\c con abr\'eg\'ee $f_{kj}f_{ji}=f_{ki}$), on
sait qu'il existe un espace fibr\'e $E$ sur~$X$, d\'efini \`a
isomorphisme pr\`es par la condition que l'on ait des isomorphismes
$f_i\colon E|X_i\isomto E_i$, satisfaisant les relations
$f_{ji}=f_jf_i^{-1}$ (avec l'abus d'\'ecriture habituel). Soit $X'$
l'espace somme des $X_i$, qui est donc un espace fibr\'e au-dessus
de~$X$ (\ie muni d'une application continue $X'\to X$). On peut
interpr\'eter de fa\c con plus concise la donn\'ee des $E_i$
comme un espace fibr\'e $E'$ sur~$X'$, et la donn\'ee des $f_{ji}$
comme un isomorphisme entre les deux images inverses (par les deux
projections canoniques) $E''_1$ et $E''_2$ de~$E'$ sur~$X''=X'\times_X
X'$, la condition de recollement pouvant alors s'\'ecrire comme une
identit\'e entre isomorphismes d'espaces fibr\'es $E'''_1$ et
$E'''_3$ sur le produit fibr\'e triple $X'''=X'\times_X X'\times_X
X'$ (o\`u $E'''_i$ d\'esigne l'image inverse de~$E'$ sur~$X'''$
par la projection canonique d'indice $i$). La construction de~$E$,
\`a partir de~$E'$ et de~$f$, est un cas typique
de proc\'ed\'e de \og descente\fg. D'ailleurs, partant d'un espace
fibr\'e $E$ sur~$X$, on dit que $X$ est \og localement trivial\fg, de
fibre $F$, s'il existe un recouvrement ouvert $(X_i)$ de~$X$ tel que
les $E|X_i$ soient isomorphes \`a $F\times X_i$, ou ce qui revient
au m\^eme, tel que l'image inverse $E'$ de~$E$ sur~$X'=\coprod_i
X_i$ soit isomorphe \`a $X'\times F$.

Ainsi, la notion de \og recollement\fg d'objets comme celle de
\og localisation\fg d'une propri\'et\'e, sont li\'ees \`a
l'\'etude de certains types de \og changements de base\fg $X'\to X$.
En G\'eom\'etrie Alg\'ebrique, bien d'autres types de changement
de base, et notamment les morphismes $X'\to X$ fid\`element plats,
doivent \^etre consid\'er\'es comme correspondant \`a un
proc\'ed\'e de \og localisation\fg relativement aux
pr\'esch\'emas, ou autres objets, \og au-dessus\fg de~$X$. Ce type de
localisation est utilis\'e tout autant que la simple localisation
topologique (qui en est un cas particulier d'ailleurs). Il en est de
m\^eme (dans une moindre mesure) en G\'eom\'etrie Analytique.

La plupart des d\'emonstrations, se r\'eduisant \`a des
v\'erifications, sont omises ou simplement esquiss\'ees; le cas
\'ech\'eant nous pr\'ecisons les diagrammes moins \'evidents
qui s'introduisent dans une d\'emonstration.

\section{Univers, cat\'egories, \'equivalence de cat\'egories}
\label{VI.1}

Pour \'eviter certaines difficult\'es logiques, nous admettrons
ici la notion d'\emph{Univers}, qui est un ensemble \og assez gros\fg
pour qu'on n'en sorte pas par les op\'erations habituelles de la
th\'eorie des ensembles; un \og \emph{axiome des Univers}\fg garantit
que tout objet se trouve dans un Univers. Pour des d\'etails, voir
un livre en pr\'eparation par C.\, Chevalley et le conf\'erencier\footnote{Les auteurs d\'efinitifs sont C.\, Chevalley et P.\, Gabriel. Le livre doit sortir en l'an 3000. En attendant, \cf aussi SGA~4 I.}. Ainsi, le sigle~$\Ens$
\label{indnot:fb}d\'esigne, non pas la cat\'egorie de tous les ensembles (notion
qui n'a pas de sens), mais la cat\'egorie des ensembles qui se
trouvent dans un Univers donn\'e (que nous ne pr\'eciserons pas
ici dans la notation). De m\^eme, $\Cat$
\label{indnot:fc}d\'esignera la cat\'egorie des cat\'egories se trouvant dans
l'Univers en question, les \og morphismes\fg d'un objet $X$ de~$\Cat$
dans un autre $Y$, \'etant par d\'efinition les \emph{foncteurs}
de~$X$ dans $Y$.

Si
$\cal{C}$ est une cat\'egorie, nous d\'esignons par $\Ob(\cal{C})$
\label{indnot:fd}\emph{l'ensemble des objets} de~$\cal{C}$, par~$\Fl(\cal{C})$
\label{indnot:fe}\emph{l'ensemble des flèches} de~$\cal{C}$ (ou morphismes de
$\cal{C}$). Nous \'ecrirons donc $X\in\Ob(\cal{C})$ en \'evitant
l'abus de notation courant $X\in\cal{C}$. Si $\cal{C}$ et $\cal{C}'$
sont deux cat\'egories, un \emph{foncteur} de~$\cal{C}$ dans
$\cal{C}'$ sera toujours ce qu'on appelle commun\'ement un foncteur
\emph{covariant} de~$\cal{C}$ dans $\cal{C}'$; sa donn\'ee implique
celle de la cat\'egorie d'arriv\'ee et la cat\'egorie de
d\'epart, $\cal{C}$ et $\cal{C}'$. Les foncteurs de~$\cal{C}$ dans
$\cal{C}'$ forment un ensemble, not\'e $\Hom(\cal{C},\cal{C}')$, qui
est l'ensemble des objets d'une cat\'egorie not\'ee
$\SheafHom(\cal{C},\cal{C}')$.
\label{indnot:ff}Par d\'efinition, un \emph{foncteur contravariant} de~$\cal{C}$ dans
$\cal{C}'$ est un foncteur de la \emph{catégorie opposée}
$\cal{C}^\circ$
\label{indnot:fg}de~$\cal{C}$ dans~$\cal{C}'$.

Nous admettrons la notion de \emph{limite projective} et de
\emph{limite inductive} d'un foncteur $F\colon \cal{I}\to \cal{C}$, et
en particulier les cas particuliers les plus courants de cette notion:
produits cart\'esiens et produits fibr\'es, notions duales de
sommes directes et de sommes amalgam\'ees, et les propri\'et\'es
formelles habituelles de ces op\'erations.

Par exemple, dans la cat\'egorie $\Cat$ introduite plus haut, les
limites projectives (relatives \`a des cat\'egories $\cal{I}$ se
trouvant dans l'Univers choisi) existent; l'ensemble d'objets
(\resp l'ensemble de fl\`eches) de la cat\'egorie limite
projective $\cal{C}$ des $\cal{C}_i$, s'obtient en prenant la limite
projective des ensembles d'objets (\resp des ensembles de fl\`eches)
des cat\'egories $\cal{C}_i$. Le cas le plus connu est celui du
produit d'une famille de cat\'egories. Nous utiliserons constamment
par la suite le produit fibr\'e de deux cat\'egories sur une
troisi\`eme.

Pour tout ce qui concerne les cat\'egories et foncteurs, en
attendant le livre en pr\'eparation d\'ej\`a signal\'e, voir
\cite{VI.1} (qui est n\'ecessairement fort incomplet, m\^eme en ce qui
concerne les g\'en\'eralit\'es esquiss\'ees dans le
pr\'esent num\'ero).

Prenons cette occasion pour expliciter la notion d'\'equivalence de
cat\'egories,
\index{equivalence@\'equivalence de cat\'egories|hyperpage}qui n'est pas expos\'ee de fa\c con satisfaisante
dans~\cite{VI.1}. Un foncteur $F\colon\cal{C}\to \cal{C}'$ est dit
\emph{fidèle}
\index{fidele (foncteur)@fid\`ele (foncteur)|hyperpage}(\resp \emph{pleinement fidèle})
\index{pleinement fid\`ele (foncteur)|hyperpage}si pour tout couple d'objets $S$, $T$ de~$\cal{C}$, l'application
$u\mto F(u)$ de~$\Hom(S,T)$ dans $\Hom(F(S),F(T))$ est injective
(\resp bijective). On dit que $F$ est une \emph{équivalence} de
cat\'egories si
$F$ est pleinement fid\`ele, et si de plus tout objet $S'$ de
$\cal{C}'$ est isomorphe \`a un objet de la forme $F(S)$. On montre
qu'il revient au m\^eme de dire qu'il existe un foncteur $G$ de
$\cal{C}'$ dans $\cal{C}$ \emph{quasi-inverse de} $F$, \iev tel que
$GF$ soit isomorphe \`a $\id_{\cal{C}}$. Lorsqu'il en est ainsi, la
donn\'ee d'un foncteur $G\colon\cal{C}'\to \cal{C}$ et d'un
isomorphisme $\varphi\colon GF\to \id_{\cal{C}'}$ \'equivaut \`a
la donn\'ee, pour tout $S'\in\Ob(\cal{C}')$, d'un couple $(S,u)$
form\'e d'un objet $S$ de~$\cal{C}$ et d'un isomorphisme $u\colon
F(S)\to S'$, soit $(G(S),\varphi(S))$. (Avec ces notations, il existe
un foncteur unique $\cal{C}'\to\cal{C}$ ayant l'application donn\'ee
$S\mto G(S)$ comme application-objets, et tel que l'application
$S\mto \varphi(S)$ soit un homomorphisme de foncteurs $FG\to
\id_{\cal{C}'}$). Enfin, si $G$ est un foncteur quasi-inverse de~$F$,
et si on choisit des isomorphismes $\varphi\colon
FG\isomto\id_{\cal{C}'}$ et $\psi\colon GF\isomto\id_{\cal{C}}$, alors
les deux conditions de compatibilit\'es sur~$\varphi$, $\psi$
\'enonc\'ees dans \cite[I.1.2]{VI.1} sont en fait \'equivalentes l'une
\`a l'autre, et pour tout isomorphisme $\varphi$ choisi, il existe
un isomorphisme $\psi$ unique tel que lesdites conditions soient
satisfaites.

\section{Cat\'egories sur une autre}
\label{VI.2}

Soit $\cal{E}$ une cat\'egorie dans~$\Univ$, c'est donc un objet
de~$\Cat$, et on peut consid\'erer la cat\'egorie
$\Cat_{/\cal{E}}$
\label{indnot:fh}des \og objets de~$\Cat$ au-dessus de~$\cal{E}$\fg. Un objet de cette
cat\'egorie est donc un foncteur
$$
p\colon \cal{F}\to\cal{E}
$$
On dit aussi que la cat\'egorie~$\cal{F}$, munie d'un tel foncteur,
est une \emph{catégorie au-dessus de}~$\cal{E}$, ou une
$\cal{E}$\emph{-catégorie}. On appellera donc $\cal{E}$-foncteur
d'une cat\'egorie $\cal{F}$ sur~$\cal{E}$ dans une cat\'egorie
$\cal{G}$ sur~$\cal{E}$, un foncteur
$$
f\colon\cal{F}\to\cal{G}
$$
tel que
$$
qf=p
$$
o\`u $p$ et $q$ sont les foncteurs-projection pour $\cal{F}$
\resp $\cal{G}$. L'ensemble des $\cal{E}$-foncteurs $f$ de~$\cal{F}$
dans $\cal{G}$ est donc en correspondance biunivoque avec l'ensemble
des fl\`eches d'origine $\cal{F}$ et d'extr\'emit\'e $\cal{G}$
dans~$\Cat_{/\cal{E}}$, sans
pourtant qu'on ait l\`a une identit\'e (puisque la donn\'ee d'un
$f$ comme dessus ne d\'etermine pas $\cal{F}$ et $\cal{G}$ en tant
que cat\'egories sur~$\cal{E}$); mais bien entendu, comme dans toute
autre cat\'egorie~$\cal{C}_{/S}$, on fera couramment l'abus de
langage consistant \`a identifier les $\cal{E}$-foncteurs (au sens
explicit\'e plus haut) \`a des fl\`eches dans une
cat\'egorie~$\Cat_{/\cal{E}}$.

On d\'esignera par
$$
\Hom_{\cal{E}}(\cal{F},\cal{G})
$$
l'ensemble des $\cal{E}$-foncteurs de~$\cal{F}$ dans~$\cal{G}$. Bien
entendu, un compos\'e de~$\cal{E}$-foncteurs est un
$\cal{E}$-foncteur (la composition en question correspondant par
d\'efinition \`a la composition des fl\`eches
dans~$\Cat_{/\cal{E}}$).

Consid\'erons maintenant deux $\cal{E}$-foncteurs
$$
f,g\colon \cal{F}\to\cal{G}
$$
et un homomorphisme de foncteurs:
$$
u\colon f\to g
$$
On dit que $u$ est un $\cal{E}$-\emph{homomorphisme} ou un
\og \emph{homomorphisme de~$\cal{E}$-foncteurs}\fg, si pour tout
$\xi\in\Ob(\cal{F})$, on a
$$
q(u(\xi))=\id_{p(\xi)} \quoi\text{,}
$$
en paroles: posant
$S=p(\xi)=qf(\xi)=qg(\xi)\in\Ob(\cal{E})$,
le morphisme
$$
u(\xi)\colon f(\xi)\to g(\xi)
$$
dans $\cal{G}$ est un $\id_S$-morphisme. (De fa\c con
g\'en\'erale, pour tout morphisme $\alpha\colon T \to S$
dans~$\cal{E}$, et toute cat\'egorie $\cal{G}$ au-dessus
de~$\cal{E}$, un morphisme $v$ dans $\cal{G}$ est appel\'e un
$\alpha$-\emph{morphisme} si $q(v)=\alpha$, $q$ d\'esignant le
foncteur projection $\cal{G}\to \cal{E}$). Si on a un troisi\`eme
$\cal{E}$-foncteur $h\colon \cal{F}\to \cal{G}$ et un
$\cal{E}$-homomorphisme $v\colon g\to h$, alors $vu$ est \'egalement
un $\cal{E}$-homomorphisme. Ainsi,
\emph{les $\cal{E}$-foncteurs de~$\cal{F}$ dans~$\cal{G}$, et
les $\cal{E}$-homomorphismes de tels, forment une sous-catégorie
de la catégorie $\SheafHom(\cal{F},\cal{G})$ de tous les foncteurs
de~$\cal{F}$ dans~$\cal{G}$, qu'on appellera la catégorie des
$\cal{E}$-foncteurs de~$\cal{F}$ dans $\cal{G}$ et qu'on notera}
$$
\SheafHom_{\cal{E}/-}(\cal{F},\cal{G})
$$
\label{indnot:fi}C'est aussi la sous-cat\'egorie noyau du couple de foncteurs
$$
\xymatrix@C=.5cm{{R,S\colon\mathbf{Hom}(\cal{F},\cal{G})}\ar@<2pt>[r]\ar@<-2pt>[r]&{\mathbf{Hom}(\cal{F},\cal{E})}}\quoi\text{,}
$$
o\`u $R$ est le foncteur constant d\'efini par l'objet $p$ de
$\SheafHom(\cal{F},\cal{E})$, et o\`u $S$ est le foncteur $f\mto
q\circ f$ d\'efini par $q\colon\cal{G}\to\cal{E}$.

Pour finir ces g\'en\'eralit\'es, il reste \`a d\'efinir les
accouplements naturels entre les cat\'egories
$\SheafHom_{\cal{E}/-}(\cal{F},\cal{G})$ par composition de
$\cal{E}$-foncteurs. En d'autres termes, on veut d\'efinir un
\og foncteur composition\fg:
\begin{equation*}
\label{eq:VI.2.i}
\tag{i} {\SheafHom_{\cal{E}/-}(\cal{F},\cal{G})\times
\SheafHom_{\cal{E}/-}(\cal{G},\cal{H})\to
\SheafHom_{\cal{E}/-}(\cal{F},\cal{H})}
\end{equation*}
lorsque $\cal{F}$, $\cal{G}$, $\cal{H}$ sont trois cat\'egories
sur~$\cal{E}$, de telle fa\c con que ce foncteur induise, pour les
objets, l'application de composition $(f,g)\mto gf$ de
$\cal{E}$-foncteurs $f\colon\cal{F}\to\cal{G}$ et $g\colon \cal{G}\to
\cal{H}$. Pour ceci, rappelons qu'on d\'efinit un foncteur
canonique:
\begin{equation*}
\label{eq:VI.2.ii}
\tag{ii}
{\SheafHom(\cal{F},\cal{G})\times\SheafHom(\cal{G},\cal{H})\to
\SheafHom(\cal{F},\cal{H})}
\end{equation*}
qui, pour les objets, n'est autre que l'application de composition
$(f,g)\mto gf$ de foncteurs, et qui transforme une fl\`eche
$(u,v)$,~o\`u
$$
u\colon f\to f'\quoi,\quad v\colon g\to g'
$$
sont des fl\`eches dans $\SheafHom(\cal{F},\cal{G})$ \resp dans
$\SheafHom(\cal{G},\cal{H})$, en la fl\`eche
$$
v*u\colon gf\to g'f'
$$
\label{indnot:fj}d\'efinie par la relation:
$$
v*u(\xi)=v(f'(\xi))\cdot g(u(\xi))=g'(u(\xi))\cdot v(f(\xi))
$$
Il est bien connu que l'on obtient bien ainsi un homomorphisme de~$gf$
dans~$g'f'$, et que (pour $f,g$ et $u,v$ variables) on obtient ainsi
un foncteur~\eqref{eq:VI.2.ii}, \ie qu'on a
\begin{equation*}
\label{eq:VI.2.I}
\tag{I} {\id_g*\id_f=\id_{gf}\quoi,}
\end{equation*}
\begin{equation*}
\label{eq:VI.2.II}
\tag{II} {(v'*u')\circ(v*u)=(v'\circ v)*(u'\circ u)}
\end{equation*}
Rappelons aussi qu'on a une formule d'associativit\'e pour les
accouplements canoniques~\eqref{eq:VI.2.ii}, qui s'exprime d'une part
par l'associativit\'e $(hg)f=h(gf)$ de la composition de foncteurs,
et d'autre part par la formule
\begin{equation*}
\label{eq:VI.2.III}
\tag{III} {(w*v)*u=w*(v*u)}
\end{equation*}
pour les produits de composition d'homomorphismes de foncteurs (o\`u
$u\colon f\to f'$ et $v\colon g\to g'$ sont comme dessus, et o\`u on
suppose donn\'e de plus un homomorphisme $w\colon h\to h'$ de
foncteurs $h,h'\colon \cal{H}\to\cal{K}$). Je dis maintenant que
\emph{lorsque~$\cal{F}$, $\cal{G}$ sont des $\cal{E}$-catégories,
le foncteur composition canonique}~\eqref{eq:VI.2.ii} \emph{induit un
foncteur}~\eqref{eq:VI.2.i}. Comme on sait d\'ej\`a que le
compos\'e de deux $\cal{E}$-foncteurs est un $\cal{E}$-foncteur,
cela revient \`a dire que \emph{lorsque $u\colon f\to f'$ et
$v\colon g\to g'$ sont des homomorphismes de~$\cal{E}$-foncteurs,
alors $v*u\colon gf\to g'f'$ est également un homomorphisme de
$\cal{E}$-foncteurs.} Cela r\'esulte en effet trivialement des
d\'efinitions. Comme les accouplements~\eqref{eq:VI.2.i} sont
induits par les accouplements~\eqref{eq:VI.2.ii}, ils satisfont \`a
la m\^eme propri\'et\'e d'associativit\'e, exprim\'ee aussi
dans les formules $(hg)f=h(gf)$ et $(w*v)*u=w*(v*u)$ pour des
$\cal{E}$-foncteurs et des $\cal{E}$-homomorphismes de
$\cal{E}$-foncteurs.

Pour compl\'eter le formulaire~\eqref{eq:VI.2.I},
\eqref{eq:VI.2.II}, \eqref{eq:VI.2.III}, rappelons aussi les formules:
\begin{equation*}
\label{eq:VI.2.IV}
\tag{IV} {v*\id_{\cal{F}}=v\quad\text{et}\quad\id_{\cal{G}}*u=u}\quoi,
\end{equation*}
o\`u
pour simplifier on \'ecrit $v*f$ ou $u*g$ au lieu de~$v*u$, lorsque
$u$ \resp $v$ est l'automorphisme identique de~$f$ \resp $g$.

De la d\'efinition des accouplements~\eqref{eq:VI.2.i} r\'esulte
que $\SheafHom_{\cal{E}/-}(\cal{F},\cal{G})$ \emph{est un foncteur en
$\cal{F}$,~$\cal{G}$,
de la catégorie produit
${\Cat_{/\cal{E}}}^\circ\times\Cat_{/\cal{E}}$ dans la
catégorie~$\Cat$}. Si en effet $g\colon \cal{G}\to\cal{G}_1$ est
un $\cal{E}$-foncteur, \ie un objet de
$\SheafHom_{\cal{E}/-}(\cal{G},\cal{G}_1)$, alors faisant
dans~\eqref{eq:VI.2.i} $\cal{H}=\cal{G}_1$, il lui correspond un
foncteur
$$
g_*\colon\SheafHom_{\cal{E}/-}(\cal{F},\cal{G})\to
\SheafHom_{\cal{E}/-}(\cal{F},\cal{G}_1)
$$
On d\'efinit de la fa\c con analogue, pour un $\cal{E}$-foncteur
$f\colon \cal{F}_1\to \cal{F}$, un foncteur
$$
f^*:\SheafHom_{\cal{E}/-}(\cal{F},\cal{G})\to
\SheafHom_{\cal{E}/-}(\cal{F}_1,\cal{G})
$$
Pour abr\'eger, on d\'esigne ces foncteurs aussi par les sigles
$f\mto g\circ f$ \resp $g\mto g\circ f$ (qui d\'esignent
seulement, en fait, les applications correspondantes sur les ensembles
d'objets). Il r\'esulte de la propri\'et\'e d'associativit\'e
indiqu\'ee plus haut que de cette fa\c con, on obtient bien comme
annonc\'e un foncteur
$\Cat_{/\cal{E}}^\circ\times\Cat_{/\cal{E}}\to\Cat$.

\section{Changement de base dans les cat\'egories sur~$\mathcal{E}$}
\label{VI.3}

Comme dans $\Cat$ les limites projectives (relativement \`a des
cat\'egories $\cal{I}$ \'el\'ements de~$\Univ$) existent, il en
est de m\^eme dans $\Cat_{/\cal{E}}$, en particulier les produits
cart\'esiens y existent, qui s'interpr\`etent comme des produits
fibr\'es dans~$\Cat$. Conform\'ement aux notations
g\'en\'erales, si $\cal{F}$ et $\cal{G}$ sont des cat\'egories
au-dessus de~$\cal{E}$, on d\'esigne par
$$
\cal{F}\times_{\cal{E}}\cal{G}
$$
\label{indnot:fk}leur produit dans~$\Cat_{/\cal{E}}$, \ie leur produit fibr\'e
au-dessus de~$\cal{E}$ dans~$\Cat$,
en tant que cat\'egorie au-dessus de~$\cal{E}$. Ainsi,
$\cal{F}\times_{\cal{E}} \cal{G}$ est muni de deux $\cal{E}$-foncteurs
$\pr_1$ et $\pr_2$, qui d\'efinissent, pour toute cat\'egorie
$\cal{H}$ au-dessus de~$\cal{E}$, une bijection
$$
\Hom_{\cal{E}}(\cal{H},\cal{F}\times_{\cal{E}}\cal{G})\isomto
\Hom_{\cal{E}}(\cal{H},\cal{F})\times\Hom_{\cal{E}}(\cal{H},\cal{G}).
$$
Cette bijection provient d'ailleurs d'un isomorphisme de
cat\'egories
$$
\SheafHom_{\cal{E}/-}(\cal{H},\cal{F}\times_{\cal{E}}
\cal{G})\isomto\SheafHom_{\cal{E}/-}(\cal{H},\cal{F})\times
\SheafHom_{\cal{E}/-}(\cal{H},\cal{G})
$$
en prenant les ensembles d'objets des deux membres, o\`u le foncteur
\'ecrit est celui qui a pour composantes les foncteurs $h\mto
\pr_1\circ h$ et $h\mto \pr_2\circ h$ du premier membre dans les deux
facteurs du second. On laisse au lecteur le soin de v\'erifier qu'on
obtient bien ainsi un isomorphisme (le fait analogue \'etant vrai,
plus g\'en\'eralement, chaque fois qu'on a une limite projective
de cat\'egories --- et non seulement dans le cas d'un produit
fibr\'e).

Rappelons d'ailleurs qu'on a (comme il a \'et\'e dit dans
le~\No \ref{VI.1}):
\begin{eqnarray*}
\Ob(\cal{F}\times_\cal{E}\cal{G})&=&
\Ob(\cal{F})\times_{\Ob(\cal{E})}\Ob(\cal{G})\\
\Fl(\cal{F}\times_\cal{E}\cal{G})&=&
\Fl(\cal{F})\times_{\Fl(\cal{E})}\Fl(\cal{G})\quoi,
\end{eqnarray*}
la composition des fl\`eches se faisant d'ailleurs composante par
composante.

Dans la suite, nous consid\'erons un foncteur
$$
\lambda\colon\cal{E}'\to\cal{E}
$$
et pour toute cat\'egorie $\cal{F}$ au-dessus de~$\cal{E}$, on
consid\`ere $\cal{F}\times_\cal{E}\cal{E}'$ comme une cat\'egorie
au-dessus de~$\cal{E}'$ gr\^ace \`a $\pr_2$; en d'autres termes,
nous interpr\'etons l'op\'eration \og produit fibr\'e\fg comme une
op\'eration \emph{\og changement de base\fg,} le foncteur
$\lambda\colon\cal{E}'\to\cal{E}$ prenant le nom de \emph{\og foncteur
de changement de base\fg.}
\index{changement de base (foncteur)|hyperpage}\index{foncteur changement de base|hyperpage}Conform\'ement aux faits g\'en\'eraux bien connus, on obtient
ainsi un foncteur, dit
\emph{foncteur changement de base} pour $\lambda$:
$$
\lambda^*\colon\Cat_{/\cal{E}}\to\Cat_{/\cal{E}'}\quoi,
$$
\label{indnot:fl}(adjoint du foncteur \og restriction de la base\fg qui, \`a toute
cat\'egorie $\cal{F}'$ au-dessus de~$\cal{E}'$, de foncteur
projection~$p'$, associe~$\cal{F}'$, consid\'er\'e comme
cat\'egorie au-dessus de~$\cal{E}$ par le foncteur~$p=\lambda
p'$). Comme il est bien connu dans le cas g\'en\'eral d'un
foncteur changement de base dans une cat\'egorie, le foncteur
changement de base \og commute aux limites projectives\fg, et en
particulier \og transforme\fg produits fibr\'es sur~$\cal{E}$ en
produits fibr\'es sur~$\cal{E}'$.

Soient $\cal{F}$ et $\cal{G}$ deux cat\'egories au-dessus
de~$\cal{E}$, on va d\'efinir un \emph{isomorphisme canonique}:
\begin{multline*}
\label{eq:VI.3.i}
\tag{i}
\SheafHom_{\cal{E}'/-}(\cal{F}',\cal{G}')\isomto
\SheafHom_{\cal{E}/-}(\cal{F}\times_{\cal{E}}\cal{E}',\cal{G})\\
\text{où }\cal{F}'=\cal{F}\times_{\cal{E}}\cal{E}'\text{, }
\cal{G}'=\cal{G}\times_{\cal{E}}\cal{E}'.
\end{multline*}
Pour ceci, consid\'erons le foncteur
$$
\pr_1\colon\cal{G}'=\cal{G}\times_{\cal{E}}\cal{E}'\to \cal{G}\quoi,
$$
et d\'efinissons~\eqref{eq:VI.3.i} par
$$
F\mto \pr_1\circ F\quoi,
$$
qui a priori d\'esigne un foncteur
\begin{equation*}
\label{eq:VI.3.ii}
\tag{ii} {\SheafHom(\cal{F}',\cal{G}')\to\SheafHom(\cal{F}',\cal{G})}
\end{equation*}
Il faut donc v\'erifier seulement que ce dernier induit un foncteur
pour les sous-cat\'egories~\eqref{eq:VI.3.i}, et que ce dernier est
un isomorphisme. Que~\eqref{eq:VI.3.ii} induise une bijection
$$
\Hom_{\cal{E}'/-}(\cal{F}',\cal{G}')\isomto
\Hom_{\cal{E}/-}(\cal{F}\times_{\cal{E}}\cal{E}',\cal{G})
$$
est la propri\'et\'e caract\'eristique du foncteur changement de
base. Il reste donc \`a
prouver que si $F$, $G$ sont des $\cal{E}'$-foncteurs
$\cal{F}'\to\cal{G}'$, alors \emph{l'application}
$$
u\mto \pr_1\circ u
$$
\emph{induit une bijection}
$$
\Hom_{\cal{E}'}(F,G)\isomto\Hom_{\cal{E}}(\pr_1\circ F,\pr_1\circ
G)\quoi.
$$
La v\'erification de ce fait est imm\'ediate, et laiss\'ee au
lecteur.

Il r\'esulte de cet isomorphisme~\eqref{eq:VI.3.i}, et de la fin du
num\'ero pr\'ec\'edent, que
$$
\SheafHom_{\cal{E}'/-}(\cal{F}\times_{\cal{E}}\cal{E}',\cal{G}
\times_{\cal{E}}\cal{E}')
$$
\emph{peut être considéré comme un foncteur en}
$\cal{E}',\cal{F},\cal{G}$, \emph{de la catégorie}
$\Cat_{/\cal{E}}^\circ\times\Cat_{/\cal{E}}^\circ\times\Cat_{/\cal{E}}$
\emph{dans la catégorie}~$\Cat$, isomorphe au foncteur d\'efini
par l'expression
$\SheafHom_{\cal{E}/-}(\cal{F}\times_{\cal{E}}\cal{E}',\cal{G})$. En
particulier, pour $\cal{F}$,~$\cal{G}$ fix\'es, on obtient un
foncteur en $\cal{E}'$, $\cal{E}'\to
\SheafHom_{\cal{E}''/-}(\cal{F}',\cal{G}')=
\SheafHom_{\cal{E}'/-}(\cal{F}\times_{\cal{E}}
\cal{E}',\cal{G}\times_{\cal{E}}\cal{E}')$, et en particulier le
$\cal{E}$-foncteur de projection $\lambda\colon \cal{E}'\to\cal{E}$
d\'efinit un morphisme \ie un foncteur
$$
\lambda^*_{\cal{F},\cal{G}}\colon
\SheafHom_{\cal{E}/-}(\cal{F},\cal{G})\to
\SheafHom_{\cal{E}'/-}(\cal{F}',\cal{G}')
$$
que nous allons expliciter. Pour les ensembles d'objets des deux
membres, c'est l'application
$$
f\mto f'=f\times_{\cal{E}}\cal{E}'
$$
qui exprime la d\'ependance fonctorielle de
$\cal{F}\times_{\cal{E}}\cal{E}'$ de l'objet $\cal{F}$
sur~$\cal{E}$. D'autre part, consid\'erons deux $\cal{E}$-foncteurs
$$
f,g\colon\cal{F}\to \cal{G}
$$
et un homomorphisme de~$\cal{E}$-foncteurs
$$
u\colon f\to g\quoi,
$$
on
va expliciter l'homomorphisme de~$\cal{E}'$-foncteurs correspondant:
$$
u'\colon f'\to g'\quoi.
$$
Pour tout
$$
\xi'=(\xi,S')\in\Ob(\cal{F}')
$$
avec
$$
\xi\in\Ob(\cal{F}),\quad S'\in\Ob(\cal{E}'),\quad p(\xi)=\lambda(S')=S
$$
le morphisme
$$
u'(\xi')\colon f'(\xi')=(f(\xi),S')\to g'(\xi')=(g(\xi),S')\quad
\text{dans } \cal{G}'
$$
est d\'efini par la formule
$$
u'(\xi')=(u(\xi),\id_{S'})
$$
(ce qui est bien un $S'$-morphisme dans~$\cal{G}'$, car
$q(u(\xi))=\lambda(\id_{S'})=\id_S$).

Consid\'erons maintenant un $\cal{E}$-foncteur quelconque
$$
\lambda'\colon\cal{E}''\to\cal{E}'
$$
et le foncteur correspondant
$$
\SheafHom_{\cal{E}'/-}(\cal{F}\times_{\cal{E}}\cal{E}',
\cal{G}\times_{\cal{E}}\cal{E}')\to
\SheafHom_{\cal{E}''/-}(\cal{F}\times_{\cal{E}}\cal{E}'',
\cal{G}\times_{\cal{E}}\cal{E}'')\quoi,
$$
je dis que ce foncteur n'est autre que le foncteur qu'on obtient par
le proc\'ed\'e pr\'ec\'edent, en partant de~$\cal{F}'$ et
$\cal{G}'$ sur~$\cal{E}'$ et en consid\'erant $\cal{E}''$ comme une
cat\'egorie sur~$\cal{E}'$, compte tenu des isomorphismes de
\emph{\og transitivité de changement de base\fg}:
$$
\cal{F}'\times_{\cal{E}'}\cal{E}''\isomto
\cal{F}''=\cal{F}\times_{\cal{E}}\cal{E}''\quad
\text{et}\quad\cal{G}'\times_{\cal{E}'}\cal{E}''\isomto\cal{G}''
=\cal G\times_{\cal E}\cal E''\quoi,
$$
qui impliquent un isomorphisme canonique
$$
\SheafHom_{\cal{E}''/-}(\cal{F}'\times_{\cal{E}'}\cal{E}'',
\cal{G}'\times_{\cal{E}'}\cal{E}'')\isomto
\SheafHom_{\cal{E}''/-}(\cal{F}\times_{\cal{E}}\cal{E}'',
\cal{G}\times_{\cal{E}}\cal{E}'')
$$
La
v\'erification de cette compatibilit\'e est imm\'ediate, et
laiss\'ee au lecteur.

Les foncteurs qu'on vient de d\'efinir sont compatibles avec les
accouplements d\'efinis au num\'ero pr\'ec\'edent, de fa\c con pr\'ecise, si $\cal{F}$, $\cal{G}$,~$\cal{H}$ sont des
cat\'egories au-dessus de~$\cal{E}$ et si on pose
$$
\cal{F}'=\cal{F}\times_{\cal{E}}\cal{E}'\quoi,\quad \cal{G}'=
\cal{G}\times_{\cal{E}}\cal{E}'\quoi,\quad \cal{H}'=
\cal{H}\times_{\cal{E}}\cal{E}'\quoi,
$$
on a commutativit\'e dans le diagramme de foncteurs suivant:
$$
\xymatrix{{\mathbf{Hom}_{\cal{E}/-}(\cal{F},\cal{G})\times
\mathbf{Hom}_{\cal{E}/-}(\cal{G},\cal{H})}\ar[r]\ar[d]_{\lambda^*_{\cal{F},\cal{G}}\times \lambda^*_{\cal{G},\cal{H}}}&{\mathbf{Hom}_{\cal{E}/-}(\cal{F},\cal{H})}\ar[d]^{\lambda^*_{\cal{F},\cal{H}}}\\{\mathbf{Hom}_{\cal{E}'/-}(\cal{F}',\cal{G}')\times
\mathbf{Hom}_{\cal{E}'/-}(\cal{G}',\cal{H}')}\ar[r]&{\mathbf{Hom}_{\cal{E}'/-}(\cal{F}',\cal{H}')}}
$$
o\`u les fl\`eches horizontales sont les foncteurs-composition
d\'efinis au num\'ero pr\'ec\'edent. Cette commutativit\'e
s'exprime par les formules
$$
(gf)'=g'f'
$$
pour $f\in\Hom_{\cal{E}}(\cal{F},\cal{G})$,
$g\in\Hom_{\cal{E}}(\cal{G},\cal{H})$, (formule qui exprime simplement
la fonctorialit\'e du changement de base), et la formule
$$
(v*u)'=v'*u'
$$
lorsque $u\colon f\to f_1$ est une fl\`eche de
$\SheafHom_{\cal{E}/-}(\cal{F},\cal{G})$ et $v\colon g\to g_1$ une
fl\`eche de~$\SheafHom_{\cal{E}/-}(\cal{G},\cal{H})$. La
v\'erification de cette formule r\'esulte facilement des
d\'efinitions.

Dans la suite, nous nous int\'eresserons surtout \`a
$\SheafHom_{\cal{E}}(\cal{F},\cal{G})$ (et certaines
sous-cat\'egories remarquables de celle-ci) lorsque
$\cal{F}=\cal{E}$, et introduisons pour cette raisons une notation
sp\'eciale:
$$
\bf{\Gamma}(\cal{G}/\cal{E})=
\SheafHom_{\cal{E}}(\cal{E},\cal{G})\quoi,\quad
\Gamma(\cal{G}/\cal{E})=\Ob(\bf{\Gamma}(\cal{G}/\cal{E}))=
\Hom_{\cal{E}}(\cal{E},\cal{G})\quoi.
$$
\label{indnot:fm}
\begin{remarquesstar}
Lorsque
$\cal{E}$ est une cat\'egorie ponctuelle,
\ie $\Ob(\cal{E})$ et $\Fl(\cal{E})$ r\'eduits \`a un seul
\'el\'ement, ce qui signifie aussi que $\cal{E}$ est un objet
final de la cat\'egorie~$\Cat$, alors la donn\'ee d'une
cat\'egorie sur~$\cal{E}$ est \'equivalente \`a la donn\'ee
d'une cat\'egorie tout court, (car il y aura un foncteur unique de
$\cal{F}$ dans~$\cal{E}$). De fa\c con plus pr\'ecise,
$\Cat_{/\cal{E}}$ est alors isomorphe \`a~$\Cat$. De plus, les
cat\'egories $\SheafHom_{\cal{E}/-}(\cal{F},\cal{G})$ ne sont alors
autres que les $\SheafHom(\cal{F},\cal{G})$. Rappelons alors que la
formule fondamentale
$$
\Hom(\cal{H},\SheafHom(\cal{F},\cal{G}))\isomto
\Hom(\cal{F}\times\cal{H},\cal{G})
$$
(isomorphisme fonctoriel en les trois arguments qui y figurent),
permet d'interpr\'eter $\SheafHom(\cal{F},\cal{G})$ axiomatiquement,
en termes internes \`a la cat\'egorie~$\Cat$, de sorte que le
formulaire connu des cat\'egories $\SheafHom$ appara\^it comme un
cas particulier d'un formulaire valable dans les cat\'egories telles
que~$\Cat$, o\`u des \og objets $\SheafHom$\fg (d\'efinis par la
formule pr\'ec\'edente) existent. Il y a une interpr\'etation
analogue de~$\SheafHom_{\cal{E}/-}(\cal{F},\cal{G})$ lorsqu'on suppose
\`a nouveau $\cal{E}$ quelconque, par la formule
$$
\Hom(\cal{H},\SheafHom_{\cal{E}/-}(\cal{F},\cal{G}))\isomto
\Hom_{\cal{E}}(\cal{F}\times\cal{H},\cal{G})
$$
(isomorphisme fonctoriel en les trois arguments). De cette fa\c con,
les propri\'et\'es formelles expos\'ees dans les \No \ref{VI.2},~\ref{VI.3} sont
des cas particuliers de r\'esultats plus g\'en\'eraux, valables
dans les cat\'egories o\`u les objets
$\SheafHom_{\cal{E}/-}(\cal{F},\cal{G})$ (lorsque $\cal{F}$,~$\cal{G}$
sont deux objets de la cat\'egorie au-dessus d'un troisi\`eme~$\cal{E}$) existent.
\end{remarquesstar}

\section[Cat\'egories-fibres]{Cat\'egories-fibres; \'equivalence de
$\mathcal{E}$-cat\'egories}
\label{VI.4}

Soit~$\cal{F}$ une cat\'egorie sur~$\cal{E}$, et soit $S \in
\Ob({\cal{E}})$. On appelle \emph{catégorie-fibre de}~$\cal{F}$
en~$S$
\index{cat\'egorie-fibre|hyperpage}la sous-cat\'egorie~$\cal{F}_S$
\label{indnot:fn}de~$\cal{F}$ image r\'eciproque de la
sous-cat\'egorie ponctuelle de~$\cal{E}$ d\'efinie par $S$. Donc
les objets de~$\cal{F}_S$ sont les objets~$\xi$ de~$\cal{F}$ tels que
$p(\xi)=S$, ses morphismes sont les morphismes~$u$ de~$\cal{F}$ tels
que $p(u)=\id_S$, \ie les $S$-morphismes dans $\cal{F}$. Bien
entendu,~$\cal{F}_S$ est canoniquement isomorphe au produit fibr\'e
$\cal{F} \times_{\cal{E}} \{S\}$, o\`u~$\{S\}$ d\'esigne la
sous-cat\'egorie ponctuelle de~$\cal{E}$ d\'efinie par~$S$, munie
de son foncteur d'inclusion dans~$\cal{E}$. Il en r\'esulte (compte
tenu de la transitivit\'e du changement de base) que si on fait le
changement de base $\lambda\colon \cal{E}' \to \cal{E}$, alors pour
tout $S' \in \Ob(\cal{E}')$, \emph{la projection $\pr_1\colon
\cal{F}' = \cal{F} \times_{\cal{E}} \cal{E}' \to \cal{F}$ induit un
isomorphisme}
$$
\cal{F}'_{S'} \to \cal{F}_S \qquad (\text{où $S=\lambda(S')$}).
$$

\begin{proposition}
\label{VI.4.1} Soit $f\colon \cal{F} \to \cal{G}$ un
$\cal{E}$-foncteur. Si~$f$ est pleinement fid\`ele, alors pour tout
changement de base $\cal{E}' \to \cal{E}$, le foncteur correspondant
$f' \colon \cal{F}' = \cal{F} \times_{\cal{E}} \cal{E}' \to \cal{G}' =
\cal{G} \times_{\cal{E}} \cal{E}'$ est pleinement fid\`ele.
\end{proposition}

La v\'erification est imm\'ediate; plus g\'en\'eralement, on
peut montrer que toute limite projective de
foncteurs
pleinement
fid\`eles
(ici,~$f$ et les foncteurs identiques
dans~$\cal{E},\cal{E}'$) est un foncteur pleinement fid\`ele.

On notera que l'assertion analogue \`a 4.1,
o\`u
\og pleinement fid\`ele\fg est remplac\'ee par \og \'equi\-va\-lence
de cat\'egories\fg, est fausse, d\'ej\`a pour~$\cal{G}=\cal{E}$.
Cependant:

\begin{proposition}
\label{VI.4.2} Soit $f \colon \cal{F} \to \cal{G}$ un
$\cal{E}$-foncteur. Les conditions suivantes sont \'equivalentes:
\begin{enumerate}
\item[(i)] Il existe un $\cal{E}$-foncteur $g \colon \cal{G} \to
\cal{F}$ et des $\cal{E}$-isomorphismes
$$
gf \isomto \id_{\cal{F}}, \quad fg \isomto \id_{\cal{G}}
.
$$
\item[(ii)] Pout toute cat\'egorie $\cal{E}'$ sur~$\cal{E}$, le
foncteur
$$
f' = f \times_{\cal{E}} \cal{E}' \colon \cal{F}' = \cal{F}
\times_{\cal{E}} \cal{E}' \to \cal{G}'=\cal{G} \times_{\cal{E}}
\cal{E}'
$$
est une \'equivalence de cat\'egories.
\item[(iii)]
$f$ est une \'equivalence de cat\'egories, et pout tout $S \in
\Ob(\cal{E})$, le foncteur $f_S \colon \cal{F}_S \to \cal{G}_S$ induit
par~$f$ est une \'equivalence de cat\'egories.
\item[(iii~bis)] $f$ est pleinement fid\`ele, et pour tout $S \in
\Ob(\cal{E})$ et tout
$\eta\in\Ob(\cal{G}_S)$,
il existe un $\xi \in \Ob(\cal{F}_S)$ et un $S$-isomorphisme
$u \colon f(\xi) \to \eta$.
\end{enumerate}
\end{proposition}

\subsubsection*{D\'emonstration} \'Evidemment (i) implique que~$f$ est une
\'equivalence de cat\'egories (notion qui se d\'efinit par la
m\^eme condition, mais o\`u les isomorphismes de foncteurs ne sont
pas astreints \`a \^etre des $\cal{E}$-morphismes). D'autre part,
il r\'esulte des fonctorialit\'es du num\'ero pr\'ec\'edent
que la condition (i) est conserv\'ee apr\`es changement de
base~$\cal{E}'\to\cal{E}$. Il s'ensuit que (i)$\To$(ii).
\'Evidemment (ii)$\To$(iii), car il suffit de faire $\cal{E}'
= \cal{E}$ et~$\cal{E}'=\{S\}$. Il est encore plus trivial que
(iii)$\To$(iii~bis), reste \`a prouver que (iii~bis)$\To$(i). Pour ceci, choisissons pour tout $\eta \in
\Ob(\cal{G})$ un $g(\eta) \in \Ob(\cal{F})$ et un isomorphisme
$u(\eta)\colon f(g(\eta))\to\eta$ qui soit tel que~$q(u(\eta)) =
\id_S$, o\`u~$S=q(\eta)$. C'est possible gr\^ace \`a la
deuxi\`eme condition (iii~bis). Comme il est connu et imm\'ediat,
le fait que~$f$ est pleinement fid\`ele implique que~$g$ peut de
fa\c con unique \^etre consid\'er\'e comme un foncteur
de~$\cal{G}$ dans~$\cal{F}$, de fa\c con que les~$u(\eta)$
d\'efinissent un homomorphisme (donc un isomorphisme) fonctoriel
$u\colon fg \isomto \id_{\cal{G}}$. De plus, par construction~$g$ est
un $\cal{E}$-foncteur et~$u$ un $\cal{E}$-homomorphisme. Aux
donn\'ees pr\'ec\'edentes correspond alors un isomorphisme
fonctoriel $v\colon gf \to \id_{\cal{F}}$, d\'efini par la condition
que $f*v=u*f$, et on constate tout de suite que c'est \'egalement un
$\cal{E}$-morphisme, cqfd.
\begin{definition}
\label{VI.4.3} Si les conditions pr\'ec\'edentes sont
v\'erifi\'ees, on dit que~$f$ est une \'equivalence de
cat\'egories sur~$\cal{E}$, ou une $\cal{E}$-\'equivalence.
\end{definition}

\begin{corollaire}
\label{VI.4.4} Supposons que le foncteur projection $p\colon
\cal{F}\to\cal{E}$ soit un foncteur transportable, \ie que pour tout
isomorphisme $\alpha\colon T \to S$ dans~$\cal{E}$ et tout objet~$\xi$
dans~$\cal{F}_T$, il existe un isomorphisme~$u$ dans~$\cal{F}$ de
source~$\xi$ tel que~$p(u)=\alpha$. Alors tout $\cal{E}$-foncteur
$f\colon \cal{F}\to\cal{G}$ qui est une \'equivalence de
cat\'egories, est une $\cal{E}$-\'equivalence.
\end{corollaire}

R\'esulte
du crit\`ere (iii~bis).

\begin{corollaire}
\label{VI.4.5} Soit $f\colon \cal{F}\to\cal{G}$ une
$\cal{E}$-\'equivalence. Alors pour toute cat\'egorie~$\cal{H}$
sur~$\cal{E}$, les foncteurs correspondants:
\begin{align*}
\SheafHom_{\cal{E}/-}(\cal{G},\cal{H}) & \to
\SheafHom_{\cal{E}/-}(\cal{F},\cal{H}) \\
\SheafHom_{\cal{E}/-}(\cal{H},\cal{F})
& \to \SheafHom_{\cal{E}/-}(\cal{H},\cal{G})
\end{align*}
(\cf \No \ref{VI.2}) sont des \'equivalences de cat\'egories.
\end{corollaire}

Cela r\'esulte du crit\`ere (i) par le raisonnement habituel.

\section{Morphismes cart\'esiens, images inverses, foncteurs
car\-t\'e\-siens}
\label{VI.5}

Soit~$\cal{F}$ une cat\'egorie sur~$\cal{E}$, de
foncteur-projection~$p$.
\enlargethispage{.5cm}

\begin{definition}
\label{VI.5.1} Consid\'erons un morphisme
$$
\alpha\colon \eta\to\xi
$$
dans~$\cal{F}$, et soient
$$
S = p(\xi), \quad T = p(\eta), \quad f=p(\alpha).
$$
On dit que~$\alpha$ est un \emph{morphisme cartésien}
\index{cart\'esien (morphisme)|hyperpage}si pour tout $\eta' \in \Ob(\cal{F}_T)$ et tout $f$-morphisme $u\colon
\eta'\to\xi$, il existe un $T$-morphisme unique $\overline{u}\colon
\eta'\to\eta$ tel que~$u=\alpha\circ\overline{u}$.
\end{definition}

Cela signifie donc que pour tout $\eta'\in\Ob(\cal{F}_T)$,
l'application $v \mto \alpha\circ v$:
\begin{equation}
\tag{i} \Hom_T(\eta',\eta) \to \Hom_f(\eta',\xi)
\end{equation}
est bijective. Cela signifie aussi que le couple $(\eta,\alpha)$
\emph{représente le foncteur en~$\eta'$} $\cal{F}^{\circ}_T \to
\Ens$ du deuxi\`eme membre. Si pour un morphisme $f\colon T \to S$
dans~$\cal{E}$ donn\'e, et un $\xi \in \Ob(\cal{F}_S)$ donn\'e, il
existe un tel couple $(\eta,\alpha)$, \ie un morphisme
cart\'esien~$\alpha$ dans~$\cal{F}$ de but~$\xi$, tel que
$p(\alpha)=f$, alors~$\eta$ est d\'etermin\'e dans $\cal{F}_T$
\`a isomorphisme unique pr\`es. On dit alors que \emph{l'image
inverse de~$\xi$ par~$f$ existe}, et un objet~$\eta$ de~$\cal{F}_T$
muni d'un $f$-morphisme cart\'esien $\alpha\colon \eta \to \xi$ est
appel\'e \emph{une image inverse de~$\xi$ par~$f$}.
\index{image inverse|hyperpage}Souvent on suppose
choisie
une telle image inverse chaque fois qu'elle existe ($\cal{F}$
\'etant fix\'e); on notera alors l'image inverse par des symboles
tels que~$f^*_{\cal{F}}(\xi)$, ou simplement~$f^*(\xi)$
\label{indnot:fo}ou~$\xi \times_S T$
lorsque ces notations n'entra\^inent pas
de confusions;
le morphisme canonique $\alpha\colon \eta \to \xi$ sera alors
not\'e, dans ce qui suit, par~$\alpha_f(\xi)$.
\label{indnot:fp}Si pour tout $\xi \in \Ob(\cal{F}_S)$, l'image inverse de~$\xi$
par~$f$ existe, on dira aussi que \emph{le foncteur image inverse
par~$f$ dans~$\cal{F}$ existe}, et~$f^*(\xi)$ devient alors un
\emph{foncteur covariant en~$\xi$}, de~$\cal{F}_S$ dans~$\cal{F}_T$.
Ceci provient du fait que le deuxi\`eme membre dans (i) d\'epend
de fa\c con covariante de~$\xi$, \ie de fa\c con pr\'ecise
d\'esigne un foncteur de~$\cal{F}^{\circ}_T \times \cal{F}_S$
dans~$\Ens$. Cette d\'ependance fonctorielle pour~$f^*(\xi)$
s'explicite ainsi: consid\'erons des $f$-morphismes cart\'esiens
$$
\alpha\colon \eta \to \xi, \quad \alpha'\colon \eta' \to \xi'
$$
et un $S$-morphisme $\lambda\colon \xi \to \xi'$, alors il existe un
$T$-morphisme et un seul $\mu\colon \eta\to\eta'$ tel que l'on ait
$$
\alpha' \mu = \lambda \alpha
$$
(comme il r\'esulte du fait que~$\alpha'$ est cart\'esien).

Notons aussi le fait imm\'ediat suivant: consid\'erons un
diagramme commutatif
$$
\xymatrix@C=1.2cm{{\xi}\ar[d]_-{\lambda}&{\eta}\ar[l]_-{\alpha}\ar[d]^-{\mu}\\{\xi'}&{\eta'}\ar[l]_-{\alpha'}}
$$
dans~$\cal{F}$,
o\`u~$\alpha$ et~$\alpha'$ sont des $f$-morphismes, et~$\lambda$ un
$S$-isomorphisme,~$\mu$ un $T$-isomorphisme. \emph{Pour que~$\alpha$
soit cartésien, il faut et il suffit que~$\alpha'$ le
soit}.

\begin{definition}
\label{VI.5.2} Un $\cal{E}$-foncteur $F\colon \cal{F}\to\cal{G}$ est
appel\'e un \emph{foncteur cartésien}
\index{cart\'esien (foncteur)|hyperpage}s'il transforme morphismes cart\'esiens en morphismes
cart\'esiens. On d\'esigne par $\SheafHom_\cart(\cal{F},\cal{G})$
\label{indnot:fq}la sous-cat\'egorie pleine de
$\SheafHom_{\cal{E}/-}(\cal{F},\cal{G})$ form\'ee des foncteurs
cart\'esiens.

Par exemple, consid\'erant~$\cal{E}$ comme une cat\'egorie
sur~$\cal{E}$ gr\^ace au foncteur identique, tout morphisme
de~$\cal{E}$ est cart\'esien, donc un foncteur cart\'esien
de~$\cal{E}$ dans~$\cal{F}$ est un foncteur section $F\colon
\cal{E}\to\cal{F}$ qui transforme tout morphisme de~$\cal{E}$ en un
morphisme cart\'esien; un tel foncteur s'appelle une \emph{section
cartésienne} de~$\cal{F}$ sur~$\cal{E}$.
\end{definition}

\begin{proposition}
\label{VI.5.3} \textup{(i)}~Un foncteur $F\colon \cal{F} \to \cal{G}$ qui est une
$\cal{E}$-\'equivalence, est un foncteur cart\'esien.
\textup{(ii)}~Soient~$F,G$ deux $\cal{E}$-foncteurs
\emph{isomorphes}~$\cal{F}\to\cal{G}$. Si l'un est cart\'esien,
l'autre l'est. \textup{(iii)}~Le compos\'e de deux foncteurs cart\'esiens
$\cal{F}\to\cal{G}$ et $\cal{G}\to\cal{H}$ est un foncteur
cart\'esien.
\end{proposition}

L'assertion (iii) est triviale sur la d\'efinition, (ii)
r\'esulte de la remarque pr\'ec\'edant~\ref{VI.5.2}, (i) r\'esulte
facilement de la d\'efinition et du crit\`ere~\ref{VI.4.2}~(iii); plus
pr\'ecis\'ement, un morphisme~$\alpha$ dans~$\cal{F}$ est
cart\'esien si et seulement si~$F(\alpha)$ l'est.

\begin{corollaire}
\label{VI.5.4}
Soit $F \colon \cal{F} \to \cal{G}$ une
$\cal{E}$-\'equivalence. Alors pour toute cat\'egorie~$\cal{H}$
sur~$\cal{E}$, les foncteurs correspondants $G \mto G \circ F$ et
$G \mto F \circ G$ induisent des \'equivalences de
cat\'egories:
\begin{eqnarray*}
\SheafHom_\cart(\cal{G},\cal{H})\lto{\approx}
\SheafHom_\cart(\cal{F},\cal{H}) \\
\SheafHom_\cart(\cal{H},\cal{F})\lto{\approx}
\SheafHom_\cart(\cal{H}, \cal{G})
\end{eqnarray*}
\end{corollaire}

Cela se d\'eduit de la fa\c con habituelle de~\ref{VI.4.2}
crit\`ere~(i) et de~\ref{VI.5.3} (i) (ii)~(iii).
On peut pr\'eciser que \emph{le $\cal{E}$-foncteur
$G\colon \cal{G} \to \cal{H}$ est cartésien si et seulement si $G
\circ F$ l'est}, et de m\^eme \emph{un $\cal{E}$-foncteur $G\colon
\cal{H} \to \cal{F}$ est cartésien si et seulement si $F \circ G$
l'est}.

Il r\'esulte de~\ref{VI.5.4}~(iii) que si on consid\`ere la
sous-cat\'egorie~$\Cat^\cart_{/\cal{E}}$
\label{indnot:fr}de~$\Cat_{/\cal{E}}$ dont les objets sont les m\^emes que ceux de
$\Cat_{/\cal{E}}$, et dont les morphismes sont les foncteurs
\emph{cartésiens} alors on a comme au~\No \ref{VI.2} des accouplements:
$$
\SheafHom_\cart(\cal{F},\cal{G})\times
\SheafHom_\cart(\cal{G},\cal{H}) \to \SheafHom_\cart(\cal{F},\cal{H})
$$
induits
par ceux du~\No \ref{VI.2}, permettant de consid\'erer
$\SheafHom_\cart(\cal{F},\cal{G})$ comme un foncteur en $\cal{F}$,
$\cal{G}$, de la cat\'egorie $\left(\Cat^\cart_{/\cal{E}}
\right)^\circ \times \Cat^\cart_{/\cal{E}}$
dans~$\Cat$.
Nous aurons besoin de cette remarque surtout pour le cas o\`u
$\cal{F} = \cal{G}$:

\begin{definition}
\label{VI.5.5}
Soit $\cal{F}$ une cat\'egorie sur~$\cal{E}$. On d\'esigne par
$$
\varprojLim\cal{F}/\cal{E}
$$
\label{indnot:fs}la cat\'egorie des $\cal{E}$-foncteurs cart\'esiens
$\cal{E}\to\cal{F}$, \ie des sections cart\'esiennes de~$\cal{F}$
sur~$\cal{E}$.
\end{definition}

D'apr\`es ce qu'on vient de dire,
$\varprojLim\cal{F}/\cal{E}$ est un foncteur
en~$\cal{F}$, de la cat\'egorie $\Cat^\cart_{/\cal{E}}$ dans la
cat\'egorie $\Cat$.

Nous verrons plus bas les relations entre cette op\'eration
$\varprojLim$ et la notion de limite projective de
cat\'egories, ainsi que de nombreux exemples.

\section[Cat\'egories fibr\'ees et cat\'egories
pr\'efibr\'ees]{Cat\'egories fibr\'ees et cat\'egories
pr\'efibr\'ees. Produits et changement de base dans icelles}
\label{VI.6}

\begin{definition}
\label{VI.6.1}
Une cat\'egorie $\cal{F}$ sur~$\cal{E}$ est appel\'ee une
\emph{catégorie fibrée}
\index{fibree (categorie)@fibr\'ee (cat\'egorie)|hyperpage}(et on dit alors que le foncteur $\cal{F}\to\cal{E}$ est
\emph{fibrant})
\index{fibrant (foncteur)|hyperpage}si elle satisfait les deux axiomes suivants:
\begin{itemize}
\item{$\Fib_\rm{I}$}
Pour tout morphisme $f\colon T\to S$ dans~$\cal{E}$, le foncteur image
inverse par $f$ dans $\cal{F}$ existe.
\item{$\Fib_\rm{II}$} Le compos\'e de deux morphismes cart\'esiens
est cart\'esien.
\end{itemize}
Une cat\'egorie $\cal{F}$ sur~$\cal{E}$ satisfaisant la condition
$\Fib_\rm{I}$ est appel\'ee une \emph{catégorie
préfibrée sur~$\cal{E}$}.
\index{prefibree (categorie)@pr\'efibr\'ee (cat\'egorie)|hyperpage}\end{definition}

Si~$\cal{F}$ est une cat\'egorie fibr\'ee
(\resp pr\'efibr\'ee) sur~$\cal{E}$, une sous-cat\'egorie
$\cal{G}$ de~$\cal{F}$ est appel\'ee une \emph{sous-catégorie
fibrée}
\index{sous-cat\'egorie fibr\'ee (\resp pr\'efibr\'ee)|hyperpage}(\resp une \emph{sous-catégorie préfibrée}) si c'est une
cat\'egorie fibr\'ee (\resp pr\'efibr\'ee) sur~$\cal{E}$, et
si de plus le foncteur d'inclusion est cart\'esien. Si par exemple
$\cal{G}$ est une sous-cat\'egorie \emph{pleine} de~$\cal{F}$, on
voit que cela signifie que pour tout morphisme $f\colon T\to S$ dans
$\cal{E}$ et pour tout $\xi \in \Ob(\cal{G}_S)$, $f^*_{\cal{F}}(\xi)$
est $T$-isomorphe \`a un objet de~$\cal{G}_T$. Un autre cas
int\'eressant est le suivant: $\cal{F}$ \'etant une cat\'egorie
fibr\'ee sur~$\cal{E}$, consid\'erons la sous-cat\'egorie
$\cal{G}$ de~$\cal{F}$ ayant m\^emes objets, et dont les morphismes
sont les morphismes \emph{cartésiens} de~$\cal{F}$; en particulier
les morphismes de~$\cal{G}_S$ sont les isomorphismes de
$\cal{F}_S$. On voit de suite que c'est bien une sous-cat\'egorie
fibr\'ee de~$\cal{F}$, car dans la bijection
$$
\Hom_T(\eta',\eta) \isomto \Hom_f(\eta',\xi)
$$
relative \`a un $f$-morphisme cart\'esien $\alpha$ dans $\cal{F}$,
aux $T$-isomorphismes du premier membre correspondent les morphismes
cart\'esiens du second. Par d\'efinition, les sections
cart\'esiennes $\cal{E}\to\cal{F}$
correspondent
alors biunivoquement aux $\cal{E}$-foncteurs quelconques
$\cal{E}\to\cal{G}$ (mais on notera que le foncteur naturel
$$
\SheafHom_{\cal{E}/-}(\cal{E},\cal{G})\to
\SheafHom_\cart(\cal{E},\cal{F})=
\varprojLim(\cal{F}/\cal{E})
$$
est fid\`ele, mais en g\'en\'eral n'est pas pleinement
fid\`ele, \ie n'est pas un isomorphisme).

\begin{remarquesstar}
Soit $\cal{F}$ une cat\'egorie sur~$\cal{E}$. Les conditions
suivantes sont \'equivalentes:~(i)~Tous les morphismes de~$\cal{F}$
sont cart\'esiens (ii)~$\cal{F}$~est une cat\'egorie fibr\'ee
sur~$\cal{E}$, et les $\cal{F}_S$ sont des groupo\"ides
(\ie tout morphisme dans $\cal{F}_S$ est un isomorphisme). On dit
alors que $\cal{F}$ est une cat\'egorie \emph{fibrée en
groupoïdes}
\index{fibree en groupoides@fibr\'ee en groupo\"ides|hyperpage}
sur~$\cal{E}$. Ce sont elles qu'on rencontre surtout en \og th\'eorie
des modules\fg. Si $\cal{E}$ est un groupo\"ide, on montre que les
conditions (i) et (ii) \'equivalent aussi \`a la suivante:
(iii)~$\cal{F}$ est un groupo\"ide, et le foncteur projection
$p\colon \cal{F}\to\cal{E}$ est transportable (\cf \ref{VI.4.4}). Par
exemple, si $\cal{E}$ et $\cal{F}$ sont des groupo\"ides tels que
$\Ob(\cal{E})$ et $\Ob(\cal{F})$
soient r\'eduits \`a un point, de sorte que $\cal{E}$ et $\cal{F}$
sont d\'efinis, \`a isomorphisme pr\`es, par des groupes $E$ et
$F$, et le foncteur $p\colon \cal{F}\to\cal{E}$ est d\'efini par un
homomorphisme de groupes $p\colon F\to E$, alors $\cal{F}$ est
fibr\'e sur~$\cal{E}$ si et seulement si $p$ est surjectif, \ie si
$p$ d\'efinit une \emph{extension} du groupe $E$ par le groupe $G =
\Ker p$.
\end{remarquesstar}

\begin{proposition}
\label{VI.6.2}
Soit $F \colon \cal{F} \to \cal{G}$ une
$\cal{E}$-\'equivalence. Pour que $\cal{F}$ soit une cat\'egorie
fibr\'ee (\resp pr\'efibr\'ee) sur~$\cal{E}$, il faut et il
suffit que $\cal{G}$ le soit.
\end{proposition}

R\'esulte facilement des d\'efinitions et de la remarque
signal\'ee plus haut qu'un morphisme
$\alpha$
dans $\cal{F}$ est cart\'esien si et seulement si $F(\alpha)$ l'est.

\begin{proposition}
\label{VI.6.3}
Soient $\cal{F}_1$, $\cal{F}_2$ deux cat\'egories sur~$\cal{E}$, et
soit $\alpha = (\alpha_1,\alpha_2)$ un morphisme dans
$\cal{F}=\cal{F}_1 \times_{\cal{E}} \cal{F}_2$. Pour que $\alpha$ soit
cart\'esien, il faut et il suffit que ses composantes le soient.
\end{proposition}

Soit en effet $\xi_i$ le but et $\eta_i$ la source de~$\alpha_i$, et
soit $f\colon T\to S$ le morphisme de~$\cal{E}$ tel que $\alpha_1$ et
$\alpha_2$ soient des $f$-morphismes. Pour tout
$\eta'=(\eta'_1,\eta'_2)$ dans $\cal{F}_T$, on a un diagramme
commutatif
$$
\xymatrix{{\mathrm{Hom}_T(\eta',\eta)}\ar[r]\ar[d]&{\mathrm{Hom}_f(\eta',\xi)}\ar[d]\\{\mathrm{Hom}_T(\eta'_1,\eta_1)\times\mathrm{Hom}_T(\eta'_2,\eta_2)}\ar[r]&{\mathrm{Hom}_f(\eta'_1,\xi_1)\times \mathrm{Hom}_f(\eta'_2,\xi_2)}}
$$
o\`u les fl\`eches verticales sont des bijections. Donc si l'une
des fl\`eches horizontales est une bijection, il en est de m\^eme
de l'autre. Cela montre d\'ej\`a que si $\alpha_1$, $\alpha_2$
sont cart\'esiens (donc la deuxi\`eme fl\`eche horizontale
bijective) alors $\alpha$ l'est. La r\'eciproque se voit en faisant
dans le diagramme ci-dessus $\eta'_i=\eta_i$ d'o\`u
$\Hom_T(\eta'_i,\eta_i) \neq \emptyset$, d'abord pour $i=2$ ce qui
prouve que $\alpha_1$ est cart\'esien, puis pour $i=1$ ce qui prouve
que $\alpha_2$ est cart\'esien.

\begin{corollaire}
\label{VI.6.4}
Soit
$\cal{F} = \cal{F}_1\times_{\cal{E}}\cal{F}_2$, et soit
$F=(F_1,F_2)$ un $\cal{E}$-foncteur $\cal{G}\to\cal{F}$. Pour que
$F$
soit cart\'esien, il faut et il suffit que $F_1$ et $F_2$ le
soient. On obtient ainsi un isomorphisme de cat\'egories
$$
\SheafHom_\cart(\cal{G}, \cal{F}_1 \times_{\cal{E}} \cal{F}_2) \isomto
\SheafHom_\cart(\cal{G},\cal{F}_1)
\times\SheafHom_\cart(\cal{G},\cal{F}_2)
$$
et en particulier (faisant $\cal{G}=\cal{E}$) un isomorphisme de
cat\'egories
$$
\varprojLim(\cal{F}_1 \times_{\cal{E}}
\cal{F}_2/\cal{E}) \isomto
\varprojLim(\cal{F}_1/\cal{E}) \times
\varprojLim(\cal{F}_2/\cal{E})
$$
\end{corollaire}

\begin{corollaire}
\label{VI.6.5}
Soient $\cal{F}_1$ et $\cal{F}_2$ deux cat\'egories fibr\'ees
(\resp pr\'efibr\'ees) au-dessus de~$\cal{E}$, alors leur produit
fibr\'e $\cal{F}=\cal{F}_1 \times_{\cal{E}}\cal{F}_2$ est une
cat\'egorie fibr\'ee (\resp pr\'efibr\'ee) sur~$\cal{E}$.
\end{corollaire}

Ces r\'esultats s'\'etendent d'ailleurs au cas du produit
fibr\'e d'une famille quelconque de cat\'egories sur~$\cal{E}$.

\begin{proposition}
\label{VI.6.6}
Soient $\cal{F}$ une cat\'egorie sur~$\cal{E}$, de
foncteur-projection $p$, et soit $\lambda \colon \cal{E}' \to \cal{E}$
un foncteur, consid\'erons
$\cal{F}'=\cal{F}\times_{\cal{E}}\cal{E}'$ comme une cat\'egorie sur~$\cal{E}'$ par le foncteur-projection $p'=p \times_{\cal{E}}
\id_{\cal{E}'}$. Soit $\alpha'$ un morphisme de~$\cal{F}'$, pour que
$\alpha'$ soit un morphisme cart\'esien, il faut et il suffit que
son image $\alpha$ dans $\cal{F}$ le soit.
\end{proposition}

La d\'emonstration est imm\'ediate et laiss\'ee au lecteur.

\begin{corollaire}
\label{VI.6.7}
Pour tout foncteur cart\'esien $F\colon \cal{F} \to \cal{G}$ de
cat\'egories sur~$\cal{E}$, le foncteur
$F'=F\times_{\cal{E}}\cal{E}'$ de
$\cal{F}'=\cal{F}\times_{\cal{E}}\cal{E}'$ dans
$\cal{G}'=\cal{G}\times_{\cal{E}}\cal{E}'$ est cart\'esien.
\end{corollaire}

Par suite, le foncteur $\SheafHom_{\cal{E}}(\cal{F},\cal{G}) \to
\SheafHom_{\cal{E}'}(\cal{F}',\cal{G}')$ consid\'er\'e dans
le
\No \ref{VI.3} induit un foncteur
$$
\SheafHom_\cart(\cal{F},\cal{G}) \to
\SheafHom_\cart(\cal{F}',\cal{G}') \rm{;}
$$
en d'autres termes, pour $\cal{F}$, $\cal{G}$ fix\'es, \emph{on peut
considérer}
$$
\SheafHom_\cart(\cal{F} \times_{\cal{E}} \cal{E}',
\cal{G}\times_{\cal{E}} \cal{E}')
$$
\emph{comme
un foncteur en~$\cal{E}'$, de la catégorie
${\Cat_{/\cal{E}}}^\circ$ dans $\Cat$}. Si on laisse varier
\'egalement $\cal{F}$, $\cal{G}$, on trouve
un foncteur de la cat\'egorie ${\Cat_{/\cal{E}}}^\circ \times \big(
\Cat^\cart_{/\cal{E}} \big)^\circ \times \Cat^\cart_{/\cal{E}}$
dans~$\Cat$. Lorsqu'on tient compte de l'isomorphisme
$$
\Hom_{\cal{E}'}(\cal{F}',\cal{G}') \isomto
\Hom_\cal{E}(\cal{F}\times_\cal{E}\cal{E}',\cal{G})
$$
envisag\'e au \No \ref{VI.3}, alors les $\cal{E}'$-foncteurs cart\'esiens
de~$\cal{F}'$ dans~$\cal{G}'$ correspondent aux $\cal{E}$-foncteurs
$\cal{F} \times_\cal{E} \cal{E}' \to \cal{G}$ qui transforment tout
morphisme dont la premi\`ere projection est un morphisme
cart\'esien de~$\cal{F}$, en un morphisme cart\'esien de
$\cal{G}$. Faisant $\cal{F}=\cal{E}$, on trouve (apr\`es changement
de notation):

\begin{corollaire}
\label{VI.6.8}
$\varprojLim(\cal{F}'/\cal{E}')$ est isomorphe \`a
la sous-cat\'egorie pleine de
$\SheafHom_{\cal{E}/-}(\cal{E}',\cal{F})$ form\'ee des
$\cal{E}$-foncteurs $\cal{E}'\to\cal{F}$ qui transforment morphismes
quelconques en morphismes cart\'esiens. En particulier, si $\cal{F}$
est une cat\'egorie fibr\'ee et si $\tilde{\cal{F}}$ est la
sous-cat\'egorie de~$\cal{F}$ dont les morphismes sont les
morphismes cart\'esiens de~$\cal{F}$, alors on a une bijection
$$
\Ob \varprojLim(\cal{F}'/\cal{E}') \isomto
\Hom_{\cal{E}/-}(\cal{E}',\tilde{\cal{F}}) \rm{.}
$$
\end{corollaire}

Cela pr\'ecise la fa\c con dont l'expression
$\varprojLim(\cal{F} \times_{\cal{E}} \cal{E}' /
\cal{E}')$ doit \^etre consid\'er\'ee comme un foncteur en
$\cal{E}'$ et en $\cal{F}$, de la cat\'egorie
${\Cat_{/\cal{E}}}^\circ \times \Cat^\cart_{/\cal{E}}$ dans la
cat\'egorie $\Cat$. On verra ult\'erieurement une d\'ependance
fonctorielle plus compl\`ete par rapport \`a $\cal{E}'$, lorsque
$\cal{F}$ est astreint \`a \^etre une cat\'egorie fibr\'ee.

\begin{corollaire}
\label{VI.6.9}
Soit $\cal{F}$ une cat\'egorie fibr\'ee (\resp pr\'efibr\'ee)
sur~$\cal{E}$, alors $\cal{F}'=\cal{F}\times_{\cal{E}}\cal{E}'$ est
une cat\'egorie fibr\'ee (\resp pr\'efibr\'ee) sur~$\cal{E}'$.
\end{corollaire}

\begin{proposition}
\label{VI.6.10} Soient $\cal{F}$, $\cal{G}$ des cat\'egories
pr\'efibr\'ees sur~$\cal{E}$, $F$ un $\cal{E}$-foncteur
cart\'esien de~$\cal{F}$ dans $\cal{G}$. Pour que $F$ soit
fid\`ele,
(\resp pleinement fid\`ele, \resp une $\cal{E}$-\'equivalence)
il faut et il suffit que pour tout
$S\in \Ob(\cal{E})$,
le foncteur induit $F_S\colon \cal{F}_S\to \cal{G}_S$ soit
fid\`ele (\resp pleinement fid\`ele, \resp une \'equivalence).
\end{proposition}

D\'emonstration imm\'ediate \`a partir des d\'efinitions.

Pour finir ce num\'ero, nous donnons quelques propri\'et\'es des
cat\'egories fibr\'ees, utilisant l'axiome $\Fib_{\mathrm{II}}$.
\begin{proposition}
\label{VI.6.11} Soit $\cal{F}$ une cat\'egorie pr\'efibr\'ee
sur~$\cal{E}$. Pour que $\cal{F}$ soit fibr\'ee, il faut et il
suffit qu'elle satisfasse la condition suivante:

$\Fib_{\mathrm{II}}'$: Soit $\alpha\colon\eta\to\xi$ un morphisme
cart\'esien dans $\cal{F}$ au-dessus du morphisme $f\colon T\to S$
de~$\cal{E}$. Pour tout morphisme $g\colon U\to T$ dans $\cal{E}$, et
tout
$\zeta\in\Ob(\cal{F}_U)$,
l'application $u\mto\alpha\circ u$:
$$
\Hom_g(\zeta,\eta)\to\Hom_{fg}(\zeta,\xi)
$$
est bijective.
\end{proposition}

En d'autres termes, dans une cat\'egorie \emph{fibrée}
sur~$\cal{E}$, les diagrammes cart\'esiens sont caract\'eris\'es
par une propri\'et\'e, plus forte a priori que celle de la
d\'efinition (qu'on obtient en faisant $g=\id_T$ dans
l'\'enonc\'e qui pr\'ec\`ede).
\begin{corollaire}
\label{VI.6.12} Soient $\cal{F}$ une cat\'egorie sur~$\cal{E}$,
$\alpha$ un morphisme dans $\cal{F}$. Pour que $\alpha$ soit un
isomorphisme, il faut que $p(\alpha)=f$ soit un isomorphisme et que
$\alpha$ soit cart\'esien; la r\'eciproque est vraie si $\cal{F}$
est fibr\'ee sur~$\cal{E}$.
\end{corollaire}

En effet, si $\alpha$ est un isomorphisme il en est \'evidemment de
m\^eme de~$f=p(\alpha)$; pour tout
$\eta'\in\Ob(\cal{F}_T)$,
l'application $u\mto \alpha\circ u$
$$
\Hom(\eta',\eta)\to \Hom(\eta',\xi)
$$
est bijective; comme $f$ est un isomorphisme, on voit de suite
qu'un
\'el\'ement du premier membre est un $T$-morphisme si et seulement
si son image dans le second est un $f$-morphisme, donc on obtient
ainsi une bijection
$$
\Hom_T(\eta',\eta)\to\Hom_f(\eta',\xi)
$$
ce
qui prouve la premi\`ere assertion. R\'eciproquement, supposons
que $f$ soit un isomorphisme et que $\alpha$ satisfasse la condition
\'enonc\'ee dans $\Fib_{\mathrm{II}'}$ (ce qui signifie donc,
lorsque $\cal{F}$ est
fibr\'ee
sur~$\cal{E}$, que $\alpha$ est cart\'esien), alors on voit tout de
suite que pour tout $\zeta\in\Ob\cal{F}$, l'application $u\mto
\alpha\circ u$ de~$\Hom(\zeta,\eta)$ dans $\Hom(\zeta,\xi)$ est
bijective, donc $\alpha$ est un isomorphisme.
\begin{corollaire}
\label{VI.6.13} Soient $\alpha\colon\eta\to\xi$ et
$\beta\colon\zeta\to \eta$ deux morphismes composables dans la
cat\'egorie $\cal{F}$ fibr\'ee sur~$\cal{E}$. Si $\alpha$ est
cart\'esien alors $\beta$ l'est si et seulement si $\alpha\beta$
l'est.
\end{corollaire}

On utilise la d\'efinition des morphismes cart\'esiens sous la
forme renforc\'ee de~\ref{VI.6.11}.

\section{Cat\'egories cliv\'ees sur~$\mathcal{E}$}
\label{VI.7}

\begin{definition}
\label{VI.7.1} Soit $\cal{F}$ une cat\'egorie sur~$\cal{E}$. On
appelle \emph{clivage}
\index{clivage|hyperpage}de~$\cal{F}$ sur~$\cal{E}$ une fonction qui attache \`a tout
$f\in\Fl(\cal{E})$ un foncteur image inverse pour $f$ dans $\cal{F}$,
soit~$f^*$. Le clivage est dit \emph{normalisé}
\index{clivage normalis\'e|hyperpage}\index{normalis\'e (clivage)|hyperpage}si $f=\id_S$ implique $f^*=\id_{\cal{F}_S}$. On appelle
\emph{catégorie clivée} (\resp \emph{catégorie clivée
normalisée})
\index{cat\'egorie cliv\'ee (\resp cliv\'ee normalis\'ee)|hyperpage}une cat\'egorie $\cal{F}$ sur~$\cal{E}$ munie d'un clivage
(\resp d'un clivage normalis\'e).
\end{definition}

Il est \'evident que $\cal{F}$ admet un clivage si et seulement si
$\cal{F}$ est pr\'efibr\'ee sur~$\cal{E}$, et alors $\cal{F}$
admet un clivage normalis\'e. L'ensemble des clivages sur~$\cal{F}$
est en correspondance biunivoque avec l'ensemble des parties $K$ de
$\Fl(\cal{F})$ satisfaisant les conditions suivantes:
\begin{enumerate}
\item[a)] Les $\alpha\in K$ sont des morphismes cart\'esiens.
\item[b)] Pour tout morphisme $f\colon T\to S$ dans $\cal{E}$ et tout
$\xi\in\Ob(\cal{F}_S)$, il existe un $f$-morphisme unique dans $K$, de
but $\xi$.
\end{enumerate}

Pour que le clivage d\'efini par $K$ soit normalis\'e, il faut et
il suffit que $K$ satisfasse de plus la condition
\begin{enumerate}
\item[c)] Les morphismes identiques dans $\cal{F}$ appartiennent
\`a~$K$.
\end{enumerate}

Les
morphismes \'el\'ements de~$K$ pourront \^etre appel\'es les
\emph{\og morphismes de transport\fg}
\index{transport (morphisme de)|hyperpage}pour le clivage envisag\'e.

La notion d'isomorphisme de cat\'egories cliv\'ees sur~$\cal{E}$
est claire. Plus g\'en\'eralement, on peut d\'efinir les
morphismes de~$\cal{E}$-cat\'egories cliv\'ees comme les foncteurs
de~$\cal{E}$-cat\'egories $\cal{F}\to \cal{G}$ qui appliquent
morphismes de transport en morphismes de transport. (Ce sont en
particulier des foncteurs cart\'esiens). De cette fa\c con les
cat\'egories cliv\'ees sur~$\cal{E}$ sont les objets d'une
cat\'egorie, la \emph{catégorie des catégories clivées
sur} $\cal{E}$. Le lecteur explicitera l'existence de produits,
li\'ee au fait que si une cat\'egorie sur~$\cal{E}$ est produit de
cat\'egories $\cal{F}_i$ sur~$\cal{E}$ munies chacune d'un clivage,
alors $\cal{F}$ est muni d'un clivage naturel correspondant. On laisse
\'egalement au lecteur d'expliciter la notion de changement de base
dans les cat\'egories cliv\'ees.

Nous d\'esignerons par $\alpha_f(\xi)$ le morphisme canonique
$$
\alpha_f(\xi)\colon f^*(\xi)\to\xi.
$$
Il est, on l'a dit, fonctoriel en $\xi$, \ie on a un homomorphisme
fonctoriel
$$
\alpha_f\colon i_Tf^*\to i_S,
$$
o\`u pour tout $S\in\Ob(\cal{E})$, $i_S$ d\'esigne le foncteur
d'inclusion
$$
i_S\colon\cal{F}_S\to \cal{F}
$$

Consid\'erons maintenant des morphismes
$$
f\colon T\to S\quad\text{ et }\quad g\colon U\to T
$$
dans $\cal{E}$, et soit $\xi\in\Ob(\cal{F}_S)$, il existe alors un
unique $U$-morphisme
$$
c_{f,g}(\xi)\colon g^*f^*(\xi)\to (fg)^*(\xi)
$$
rendant
commutatif le diagramme
$$
\xymatrix{{f^*(\xi)}\ar[d]_{\alpha_f(\xi)}& &{g^*(f^*(\xi))}\ar[ll]_{\alpha_g(f^*(\xi))}\ar[d]^{c_{f,g}(\xi)}\\{\xi}& &{(fg)^*(\xi)}\ar[ll]^{\alpha_{fg}(\xi)}}
$$
(en vertu de la d\'efinition de~$(fg)^*(\xi)$). Pour $\xi$ variable,
cet homomorphisme est fonctoriel, \ie : on a un homomorphisme
$$
c_{f,g}\colon g^*f^*\to (fg)^*
$$
de foncteurs $\cal{F}_S\to\cal{F}_U$. Notons tout de suite:
\begin{proposition}
\label{VI.7.2} Pour que la cat\'egorie cliv\'ee $\cal{F}$
sur~$\cal{E}$ soit fibr\'ee, il faut et il suffit que les $c_{f,g}$
soient des isomorphismes.
\end{proposition}

On en conclut, prenant pour $f$ un isomorphisme, pour $g$ son inverse,
et en consid\'erant les isomorphismes $c_{f,g}$ et $c_{g,f}$:
\begin{corollaire}
\label{VI.7.3} Si $\cal{F}$ est une cat\'egorie \emph{fibrée}
cliv\'ee sur~$\cal{E}$, alors pour tout isomorphisme $f\colon T\to
S$ dans $\cal{E}$, $f^*$ est une \'equivalence de cat\'egories
$\cal{F}_S\to \cal{F}_T$.
\end{corollaire}

\begin{proposition}
\label{VI.7.4} Soit $\cal{F}$ une cat\'egorie cliv\'ee sur~$\cal{E}$.
On a

$\qquad\qquad A)\qquad
\left\{\begin{array}{ccc}c_{f,\id_T}(\xi)&=&\alpha_{\id_T}(f^*(\xi))\\
c_{\id_S,f}(\xi)&=&f^*(\alpha_{\id_S}(\xi))
\end{array}\right.$

\smallskip
$\qquad\qquad B)\qquad c_{f,gh}(\xi)\cdot
c_{g,h}(f^*(\xi))=c_{fg,h}(\xi)\cdot h^*(c_{f,g}(\xi))$

\noindent(Dans
ces formules, $f,g,h$ d\'esignent des morphismes
$$
V\to U\to T\to S
$$
et $\xi$ un objet de~$\cal{F}_S$).
\end{proposition}

La premi\`ere et seconde relation, dans le cas d'un clivage
normalis\'e, prennent la forme plus simple
\begin{equation*}
\label{eq:VI.7.A'}\quad A')\qquad {c_{f,\id_T}=\id_{f^*},\quad
c_{\id_S,f}=\id_{f^*}.\qquad\qquad\qquad\qquad\qquad\qquad}
\end{equation*}

Quant \`a la troisi\`eme, elle se visualise par la
commutativit\'e du diagramme
\begin{equation*}
\label{eq:VI.7.D} \tag{$D$}
\begin{split}
\xymatrix@C=2.2cm{{h^*g^*f^*(\xi)}\ar[r]^-{c_{g,h}(f^*(\xi))}\ar[d]_{h^*(c_{f,g}(\xi))}&{(gh)^*(f^*(\xi))}\ar[d]^{c_{f,gh}(\xi)}\\{h^*(fg)^*(\xi)}\ar[r]_-{c_{fg,h}(\xi)}&{(fgh)^*(\xi)}}
\end{split}
\end{equation*}
Dans le cas des cat\'egories
fibr\'ees
(o\`u les $c_{f,g}$ sont des isomorphismes), cette commutativit\'e
peut s'exprimer intuitivement par le fait que \emph{l'utilisation
successive des isomorphismes de la forme $c_{f,g}$ ne conduit pas
à des \og identifications contradictoires\fg.} On peut \'ecrire
\'egalement cette formule sans argument $\xi$, par l'utilisation du
produit de convolution d'homomorphismes de foncteurs:
$$
c_{fg,h}\circ (h^**c_{f,g})=c_{f,gh}\circ(c_{g,h}*f^*).
$$

\enlargethispage{.5cm}
La d\'emonstration des deux premi\`eres formules~\ref{VI.7.4} est
triviale, esquissons celle de la troisi\`eme. Pour ceci,
consid\'erons, en plus du carr\'e $(D)$, le carr\'e
d'homomorphismes:
\begin{equation*}
\label{eq:VI.7.D'} \tag{$D'$}
\begin{split}
\xymatrix@C=2cm{{g^*f^*(\xi)}\ar[r]^-{\alpha_g(f^*(\xi))}\ar[d]_{c_{f,g}(\xi)}&{f^*(\xi)}\ar[d]^{\alpha_f(\xi)}\\{(fg)^*(\xi)}\ar[r]_-{\alpha_{fg}(\xi)}&{\xi}}
\end{split}
\end{equation*}
qui
est commutatif par d\'efinition de
$c_{f,g}(\xi)$. Consid\'erons le diagramme obtenu en joignant les
sommets de~$(D)$ aux sommets correspondants de~$(D')$ par les
homomorphismes de la forme $\alpha$:
$$
\begin{array}{ll}
\alpha_h(g^*f^*(\xi)),&\alpha_{gh}(f^*(\xi))\\
\alpha_h((fg)^*(\xi)),&\alpha_{fgh}(\xi).
\end{array}
$$
Les quatre faces lat\'erales du cube ainsi obtenu sont \'egalement
commutatives: pour celle de gauche, cela provient du fait que la
colonne gauche de~$(D)$ se d\'eduit de la colonne gauche de~$(D')$
par application de~$h$, et que $\alpha_h$ est un homomorphisme
fonctoriel; pour les trois autres, ce n'est autre que la
d\'efinition des op\'erations $c$ des trois c\^ot\'es restants
de~$(D)$. Ainsi les cinq faces du cube autres que la face
sup\'erieure sont commutatives. Il en r\'esulte que les deux
$(fgh)$-morphismes $h^*g^*f^*(\xi)\to(fgh)^*(\xi)$ d\'efinis par
$(D)$ ont un m\^eme compos\'e avec $\alpha_{fgh}(\xi)\colon
(fgh)^*(\xi)\to\xi$, donc ils sont \'egaux par d\'efinition de
$(fgh)^*$.

Bornons-nous pour la suite aux cat\'egories cliv\'ees
\emph{normalisées}. Une telle cat\'egorie donne naissance aux
objets suivants:
\begin{enumerate}
\item[a)] Une application $S\mto \cal{F}_S$ de~$\Ob(\cal{E})$ dans
$\Cat$.
\item[b)] Une application $f\mto f^*$, associant \`a toute
$f\in\Fl(\cal{E})$, de source $T$ et de but $S$,
un
foncteur $f^*\colon \cal{F}_S\to\cal{F}_T$.
\item[c)] Une application $(f,g)\mto c_{f,g}$, associant \`a tout
couple de fl\`eches $(f,g)$ de~$\cal{E}$, un homomorphisme
fonctoriel $c_{f,g}\colon g^*f^*\to (fg)^*$.
\end{enumerate}

D'ailleurs ces donn\'ees satisfont aux conditions exprim\'ees dans
les formules $A')$ et~$B)$ donn\'ees plus haut. (N.B. Si on ne
s'\'etait pas born\'e au cas d'un clivage normalis\'e, il aurait
fallu introduire un objet suppl\'ementaire, savoir une fonction
$S\mto \alpha_S$ qui associe \`a tout objet $S$ de~$\cal{E}$ un
homomorphisme fonctoriel $\alpha_S\colon (\id_S)^*\to
\id_{\cal{F}_S}$; la condition $A')$ se remplacerait alors par la
condition $A)$).

Nous
allons montrer maintenant comment on peut reconstituer (\`a
isomorphisme unique pr\`es) la cat\'egorie cliv\'ee
normalis\'ee $\cal{F}$ sur~$\cal{E}$ \`a l'aide des objets
pr\'ec\'edents.

\section[Cat\'egorie cliv\'ee d\'efinie par un pseudo-foncteur]{Cat\'egorie cliv\'ee d\'efinie par un pseudo-foncteur
$\mathcal{E}^\circ\to\Cat$} \label{VI.8}

Appelons, pour abr\'eger, \emph{pseudo-foncteur} de~$\cal{E}^\circ$
dans $\Cat$ (il faudrait dire, pseudo-foncteur \emph{normalisé}),
un ensemble de donn\'ees a), b), c) comme ci-dessus, satisfaisant les
conditions $A')$ et $B)$. Au num\'ero pr\'ec\'edent, nous avons
associ\'e, \`a une cat\'egorie cliv\'ee normalis\'ee
sur~$\cal{E}$, un pseudo-foncteur $\cal{E}^\circ\to\Cat$, ici nous
allons indiquer la construction inverse. Nous laisserons au lecteur la
v\'erification de la plupart des d\'etails, ainsi que du fait que
ces constructions sont bien \og inverses\fg l'une de l'autre. De
fa\c con pr\'ecise, il y aurait lieu de consid\'erer les
pseudo-foncteurs $\cal{E}^\circ\to\Cat$ comme les objets d'une
nouvelle cat\'egorie, et de montrer que nos constructions
fournissent des \'equivalences, quasi-inverses l'une de l'autre,
entre cette derni\`ere et la cat\'egorie des cat\'egories
cliv\'ees au-dessus de~$\cal{E}$, d\'efinie au num\'ero
pr\'ec\'edent.

On pose
$$
\cal{F}_\circ=\underset{S\in\Ob(\cal{E})}\coprod
\Ob(\cal{F}(S)),
$$
ensemble somme des ensembles
$\Ob(\cal{F}(S))$
(N.B. nous noterons ici $\cal{F}(S)$ et non $\cal{F}_S$ la valeur en
l'objet $S$ de~$\cal{E}$ du pseudo-foncteur donn\'e, pour \'eviter
des confusions de notation par la suite). On a donc une application
\'evidente:
$$
p_\circ\colon\cal{F}_\circ\to\Ob\cal{E}.
$$
Soient
$$
\overline{\xi}
=(S,\xi),\quad \overline{\eta}=(T,\eta)\qquad\qquad
\text{(avec $\xi\in\Ob\cal{F}(S)$, $\eta\in\Ob\cal{F}(T)$)}
$$
deux \'el\'ements de~$\cal{F}_\circ$, et soit $f\in\Hom(T,S)$, on
posera
$$
h_f(\overline{\eta},\overline{\xi})=\Hom_{\cal{F}(T)}(\eta,f^*(\xi)).
$$
Si on a de plus un morphisme $g\colon U\to T$ dans $\cal{E}$, et un
$\zeta\in\Ob\cal{F}(U)$, on d\'efinit une application, not\'ee
$(u,v)\mto u\circ v$:
$$
h_f(\overline{\eta},\overline{\xi})\times
h_g(\overline{\zeta},\overline{\eta})\to
h_{fg}(\overline{\zeta},\overline{\xi}),
$$
\ie une application
$$
\Hom_{\cal{F}(T)}(\eta,f^*(\xi))\times\Hom_{\cal{F}(U)}(\zeta,g^*(\eta))
\to \Hom_{\cal{F}(U)}(\zeta,(fg)^*(\xi)),
$$
par la formule
$$
u\circ v=c_{f,g}(\xi)\cdot g^*(u)\cdot v,
$$
\ie $u\circ v$ est le compos\'e de la s\'equence
$$
\zeta\lto{u}
g^*(\eta)\lto{g^*(u)}
g^*f^*(\xi)\lto{c_{f,g}(\xi)} (fg)^*(\xi).
$$
On posera d'autre part
$$
h(\overline{\eta},\overline{\xi})=\underset{f\in\Hom(T,S)}\coprod
h_f(\overline{\eta},\overline{\xi}),
$$
et les accouplements pr\'ec\'edents d\'efinissent des
accouplements
$$
h(\overline{\eta},\overline{\xi})\times
h(\overline{\zeta},\overline{\eta})\to
h(\overline{\zeta},\overline{\xi}),
$$
tandis que la d\'efinition des $h(\overline{\eta},\overline{\xi})$
implique une application \'evidente:
$$
p_{\overline{\eta},\overline{\xi}}\colon
h(\overline{\eta},\overline{\xi})\to \Hom(T,S).
$$
Ceci dit, on v\'erifie les points suivants:
\begin{enumerate}
\item[1)] La composition entre \'el\'ements des
$h(\overline{\eta},\overline{\xi})$
est \emph{associative}.
\item[2)]
Pour tout $\overline{\xi}=(\xi,S)$ dans $\cal{F}_\circ$,
consid\'erons l'\'el\'ement
identit\'e
de
$$
h_{\id_S}(\overline{\xi},\overline{\xi})=
\Hom_{\cal{F}(S)}(\id_S^*(\xi),\xi)=\Hom_{\cal{F}(S)}(\xi,\xi),
$$
et son image dans $h(\overline{\xi},\overline{\xi})$. Cet objet est
une \emph{unité} \`a gauche et \`a droite pour la composition
entre \'el\'ements des $h(\overline{\eta},\overline{\xi})$.

Cela montre d\'ej\`a que \emph{l'on obtient une catégorie}
$\cal{F}$, en posant
$$
\Ob\cal{F}=\cal{F}_\circ,\quad \Fl{\cal{F}}=
\underset{\overline{\xi},\overline{\eta}\in\cal{F}_\circ}
\coprod h(\overline{\eta},\overline{\xi}).
$$
(N.B. on ne peut prendre simplement pour $\Fl\cal{F}$ la
\emph{réunion} des ensembles $h(\overline{\eta},\overline{\xi})$,
car ces derniers ne sont pas n\'ecessairement disjoints). De plus:
\item[3)] Les applications $p_\circ\colon \Ob\cal{F}\to\Ob\cal{E}$ et
$p_1=(p_{\overline{\eta},\overline{\xi}})\colon\Fl\cal{F}\to\Fl{\cal{E}}$
d\'efinissent un \emph{foncteur} $p\colon\cal{F}\to\cal{E}$. De
cette fa\c con, $\cal{F}$ devient une cat\'egorie sur~$\cal{E}$, de
plus l'application \'evidente
$h_f(\overline{\eta},\overline{\xi})\to\Hom(\overline{\eta},\overline{\xi})$
induit une \emph{bijection}
$$
h_f(\overline{\eta},\overline{\xi})\isomto
\Hom_f(\overline{\eta},\overline{\xi}).
$$
\item[4)] Les applications \'evidentes
$$ \Ob\cal{F}(S)\to\cal{F}_\circ=\Ob\cal{F},\qquad \Fl\cal{F}(S)\to
\Fl\cal{F},
$$
o\`u la deuxi\`eme est d\'efinie par les applications
\'evidentes
$$
\Hom_{\cal{F}(S)}(\xi,\xi')=
h_{\id_S}(\overline{\xi},\overline{\xi}')\to
\Hom(\overline{\xi},\overline{\xi}')
$$
d\'efinissent un \emph{isomorphisme}
$$
i_S \colon \cal{F}(S)\isomto\cal{F}_S.
$$
\item[5)] Pour tout objet $\overline{\xi}=(S,\xi)$ de~$\cal{F}$, et
tout morphisme $f\colon T\to S$ de~$\cal{E}$, consid\'erons
l'\'el\'ement $\overline{\eta}=(T,\eta)$ de~$\cal{F}_T$, avec
$\eta=f^*(\xi)$, et l'\'el\'ement $\alpha_f(\xi)$ de
$\Hom(\overline{\eta},\overline{\xi})$, image de~$\id_{f^*(\xi)}$ par
le morphisme
$\Hom_{\cal{F}(T)}(f^*(\xi),f^*(\xi))=h_f(\overline{\eta},\overline{\xi})
\to \Hom_f(\overline{\eta},\overline{\xi})$. \emph{Cet élément
est cartésien, et c'est l'identité dans $\overline{\xi}$ si}
$f=\id_S$, en d'autres termes, l'ensemble des $\alpha_f(\xi)$
d\'efinit un \emph{ clivage normalisé de~$\cal{F}$ sur}
$\cal{E}$. De plus, par construction, on a commutativit\'e dans le
diagramme de foncteurs
$$
\xymatrix{{\cal{F}(S)}\ar[r]^{f^*}\ar[d]_{i_S}&{\cal{F}(T)}\ar[d]^{i_T}\\{\cal{F}_S}\ar[r]^{f^*_{\cal{F}}}&{\cal{F}_T}}
$$
o\`u $f^*_{\cal{F}}$ est le foncteur image inverse par $f$, relatif
au clivage consid\'er\'e sur~$\cal{F}$. Enfin:
\item[6)] les homomorphismes $c_{f,g}$ donn\'es avec le
pseudo-foncteur sont transform\'es, par les isomorphismes $i_S$, en
les homomorphismes fonctoriels $c_{f,g}$ associ\'es au clivage
de~$\cal{F}$.
\end{enumerate}

Nous nous bornons \`a donner la v\'erification de 1) (qui est, si
possible, moins triviale que les autres). Il suffit de prouver
l'associativit\'e de la composition entre les objets d'ensembles de
la forme $h_f(\overline{\eta},\overline{\xi})$. Consid\'erons donc
dans $\cal{E}$ des morphismes
$$
S\lfrom{f} T\lfrom{g} U\lfrom{h} V
$$
et des objets
$$
\xi,\ \eta,\ \zeta,\ \tau
$$
dans $\cal{F}(S),\cal{F}(T),\cal{F}(U),\cal{F}(V)$, enfin des
\'el\'ements
$$
\begin{array}{ccc}
u\in
h_f(\overline{\eta},\overline{\xi})&=&\Hom_{\cal{F}(T)}(\eta,f^*(\xi))\\
v\in
h_g(\overline{\zeta},\overline{\eta})&=&\Hom_{\cal{F}(U)}(\zeta,g^*(\eta))\\
w\in
h_h(\overline{\tau},\overline{\zeta})&=&\Hom_{\cal{F}(V)}(\tau,h^*(\zeta)).
\end{array}
$$
On
veut prouver la formule
$$
(u\circ v)\circ w = u\circ (v\circ w),
$$
qui est une \'egalit\'e dans
$\Hom_{\cal{F}(V)}(\tau,(fgh)^*(\xi))$. En vertu des d\'efinitions
les deux membres de cette \'egalit\'e s'obtiennent par composition
suivant le contour sup\'erieur et inf\'erieur du diagramme
ci-dessous:
$$
\xymatrix{
\tau \ar[r]^w \ar[drr]_{v\circ w} & h^*(\zeta)
\ar[r]^{h^*(v)} \ar@<-2pt>
`u[rrrrr] `[rrrrr]^{h^*(u\circ v)} [rrrrr] & h^*g^*(\eta)
\ar[rr]^{h^*g^*(u)}
\ar[d]_{c_{g,h}(\eta)} & & h^*g^*f^*(\xi) \ar[rr]^{h^*(c_{f,g}(\xi))}
\ar[d]^{c_{g,h}(f^*(\xi))} & & h^*(fg)^*(\xi) \ar[d]^{c_{fg,h}(\xi)} \\
& & (gh)^*(\eta) \ar[rr]^{(gh)^*(u)} & & (gh)^*f^*(\xi)
\ar[rr]^{c_{f,gh}(\xi)} & & (fgh)^*(\xi) }
$$
Or le carr\'e m\'edian est commutatif parce que $c_{g,h}$ est un
homomorphisme fonctoriel, et le carr\'e de droite est commutatif en
vertu de la condition~$B)$ pour un pseudo-foncteur. D'o\`u le
r\'esultat annonc\'e.

Bien entendu, il reste \`a pr\'eciser, lorsque le pseudo-foncteur
envisag\'e provient d\'ej\`a d'une cat\'egorie cliv\'ee
normalis\'ee $\cal{F}'$ sur~$\cal{E}$, comment on obtient un
isomorphisme naturel entre $\cal{F}'$ et $\cal{F}$. Nous en laissons
le d\'etail au lecteur.

Nous laissons \'egalement au lecteur d'interpr\'eter, en termes de
pseudo-foncteurs, la notion d'image inverse d'une cat\'egorie
cliv\'ee $\cal{F}$ sur~$\cal{E}$ par un foncteur changement de base
$\cal{E}' \to\cal{E}$.

\section[Exemple]{Exemple: cat\'egorie cliv\'ee d\'efinie par un foncteur
$\cal{E}^\circ\to\Cat$; cat\'egories scind\'ees sur~$\cal{E}$}
\label{VI.9}
Supposons qu'on ait un foncteur
$$
\varphi\colon\cal{E}^\circ\to\Cat,
$$
il d\'efinit alors un pseudo-foncteur en posant
$$
\cal{F}(S)=\varphi(S),
\quad f^*=\varphi(f),\quad c_{f,g}=\id_{(fg)^*}
$$
Donc la construction du num\'ero pr\'ec\'edent nous donne une
cat\'egorie $\cal{F}$ cliv\'ee sur~$\cal{E}$, dite associ\'ee au
foncteur $\varphi$. Pour qu'une cat\'egorie cliv\'ee sur~$\cal{E}$
soit isomorphe \`a une cat\'egorie cliv\'ee d\'efinie par un
foncteur $\varphi:\cal{E}^\circ\to\Cat$, il faut et il suffit
manifestement qu'elle satisfasse les conditions:
$$
(fg)^*=g^*f^*,\quad c_{f,g}=\id_{(fg)^*}\quoi.
$$
En termes de l'ensemble $K$ des morphismes de transport, cela signifie
aussi simplement que \emph{le composé de deux morphismes de
transport est un morphisme de transport}. Un clivage d'une
cat\'egorie $\cal{F}$ sur~$\cal{E}$ satisfaisant la condition
pr\'ec\'edente est appel\'e un \emph{scindage}
\index{scindage|hyperpage}de~$\cal{F}$ sur~$\cal{E}$, et une cat\'egorie $\cal{F}$
sur~$\cal{E}$ munie d'un scindage est appel\'ee une
\emph{catégorie scindée sur~$\cal{E}$}.
\index{scindee (categorie)@scind\'ee (cat\'egorie)|hyperpage}C'est donc un cas particulier de la notion de cat\'egorie
cliv\'ee. La cat\'egorie des cat\'egories scind\'ees
sur~$\cal{E}$ est donc \'equivalente \`a
$\SheafHom(\cal{E}^\circ,\Cat)$. Noter qu'une cat\'egorie
scind\'ee sur~$\cal{E}$ est a fortiori une cat\'egorie cliv\'ee
sur~$\cal{E}$.

Si $\cal{F}$ est une cat\'egorie fibr\'ee sur~$\cal{E}$, il
n'existe pas toujours de scindage sur~$\cal{F}$. Supposons par exemple
que $\Ob\cal{E}$ et $\Ob\cal{F}$ soient r\'eduits \`a un
\'el\'ement, et que l'ensemble des endomorphismes dudit est un
groupe $E$ \resp $F$, de sorte que le foncteur projection $p$ est
donn\'e par un homomorphisme de groupes $p\colon F\to E$, surjectif
puisque $p$ est fibrant. On v\'erifie alors aussit\^ot que
l'ensemble des clivages de~$\cal{F}$ sur~$\cal{E}$ est en
correspondance biunivoque avec l'ensemble des applications $s\colon
E\to F$ telles que $ps=\id_E$ (\ie l'ensemble des \og syst\`emes de
repr\'esentants\fg pour les classes mod le sous-groupe $G$ noyau de
l'homomorphisme surjectif $p\colon F\to E$). Un clivage est un
scindage si et seulement si $s$ est un homomorphisme de groupes. Dire
qu'il existe un scindage signifie donc que l'extension de groupes $F$
de~$E$ par $G$ est triviale, ce qui s'exprime, lorsque $G$ est
commutatif, par la nullit\'e d'une certaine classe de cohomologie
dans $\H^2(E,G)$ (o\`u $G$ est consid\'er\'e comme un groupe
o\`u $E$ op\`ere).

Supposons cependant que $\cal{F}$ soit une cat\'egorie fibr\'ee
sur~$\cal{E}$ telle que les
$\cal{F}_S$ soient des cat\'egories \emph{rigides}, \ie le groupe
des automorphismes de tout objet de~$\cal{F}_S$ est r\'eduit \`a
l'identit\'e. Il est facile alors de prouver que $\cal{F}$ admet un
scindage sur~$\cal{E}$. En effet, on constate d'abord que la question
d'existence d'un scindage n'est pas modifi\'ee si on remplace
$\cal{F}$ par une cat\'egorie $\cal{E}$-\'equivalente, ce qui nous
ram\`ene en l'occurrence au cas o\`u les $\cal{F}_S$ sont des
cat\'egories rigides \emph{et réduites} (\ie deux objets
isomorphes dans $\cal{F}_S$ sont identiques). Mais si $G$ est une
cat\'egorie rigide et r\'eduite, tout isomorphisme de deux
foncteurs $H\to G$ (o\`u $H$ est une cat\'egorie quelconque) est
une identit\'e. Il s'ensuit que si $\cal{F}$ est une cat\'egorie
fibr\'ee sur~$\cal{E}$, telle que les cat\'egories-fibres soient
rigides et r\'eduites, alors il existe un clivage \emph{unique} de
$\cal{F}$ sur~$\cal{E}$, qui est n\'ecessairement un scindage. Donc
$\cal{F}$ est isomorphe \`a la cat\'egorie d\'efinie par un
foncteur $\varphi\colon \cal{E}^\circ\to \Cat$, tel que les
$\varphi(S)$ soient des cat\'egories rigides et discr\`etes, et le
foncteur $\varphi$ est d\'efini \`a isomorphisme pr\`es.

\section{Cat\'egories co-fibr\'ees, cat\'egories bi-fibr\'ees}
\label{VI.10}
Consid\'erons une cat\'egorie $\cal{F}$ au-dessus de~$\cal{E}$,
avec le foncteur projection
$$
p\colon \cal{F}\to \cal{E},
$$
elle d\'efinit une cat\'egorie $\cal{F}^\circ$ au-dessus de
$\cal{E}^\circ$, par le foncteur projection
$$
p^\circ\colon \cal{F}^\circ\to \cal{E}^\circ.
$$
Un morphisme $\alpha\colon \eta\to\xi$ dans $\cal{F}$ est dit
\emph{co-cartésien} si c'est un morphisme cart\'esien pour
$\cal{F}^\circ$ sur~$\cal{E}^\circ$. Explicitant, on voit que cela
signifie que pour tout objet $\xi'$ de~$\cal{F}_S$, l'application
$u\mto u\circ \alpha$
$$
\Hom_S(\xi,\xi')\to\Hom_f(\eta,\xi')
$$
est bijective. On dit alors aussi que $(\xi,\alpha)$ est une
\emph{image directe}
\index{image directe|hyperpage}de~$\eta$ par $f$, dans la cat\'egorie $\cal{F}$ sur~$\cal{E}$. Si
elle existe pour tout $\eta$ dans $\cal{F}_T$, on dit que le foncteur
image directe par $f$ existe, et on note ce foncteur
$f_*^{\cal{F}}$ ou~$f_*$,
\label{indnot:ft}une fois choisi. Il est donc d\'efini par
un isomorphisme de bifoncteurs sur~$\cal{F}_T^\circ\times\cal{F}_S$:
$$
\Hom_S(f_*(\eta),\xi)\isomto \Hom_f(\eta,\xi).
$$
Si donc $f_*$ existe, pour que $f^*$ existe, il faut et il suffit que
$f_*$
admette un foncteur adjoint, \ie qu'il existe un foncteur
$f^*\colon \cal{F}_S\to\cal{F}_T$ et un isomorphisme de bifoncteurs
$$
\Hom_S(f_*(\eta),\xi)\isomto \Hom_T(\eta,f^*(\xi)).
$$
Soit $g\colon U\to T$ un autre morphisme dans $\cal{E}$, et supposons
que les images inverses et directes par $f,g$ et $fg$
existent. Consid\'erons alors les homomorphismes fonctoriels
$$
\begin{array}{cccc}
c^{f,g}\colon&f_*g_*&\dpl\from&(fg)_*\\
c_{f,g}\colon&g^*f^*&\dpl\to&(fg)^*.
\end{array}
$$
On constate que si on consid\`ere $f_*g_*$ et $g^*f^*$ comme un
couple de foncteurs adjoints, ainsi que $(fg)_*$ et $(fg)^*$, les deux
homomorphismes pr\'ec\'edents sont adjoints l'un de l'autre. Donc
l'un est un isomorphisme si et seulement si l'autre l'est. En
particulier:
\begin{proposition}
\label{VI.10.1} Supposons que la cat\'egorie $\cal{F}$ sur~$\cal{E}$
soit pr\'efibr\'ee et co-pr\'efibr\'ee. Pour qu'elle soit
fibr\'ee, il faut et il suffit qu'elle soit
co-fibr\'ee.
\end{proposition}

Bien entendu, on dit que $\cal{F}$ est co-pr\'efibr\'ee \resp co-fibr\'ee
\index{co-fibr\'ee (cat\'egorie)|hyperpage}sur~$\cal{E}$, si $\cal{F}^\circ$ est pr\'efibr\'ee \resp fibr\'ee sur~$\cal{E}$. Nous dirons que $\cal{F}$ est bi-fibr\'ee
sur~$\cal{E}$,
\index{bi-fibr\'ee (cat\'egorie)|hyperpage}si elle est \`a la fois fibr\'ee et co-fibr\'ee sur~$\cal{E}$.

\section{Exemples divers}
\label{VI.11}

\begin{enumerate}

\item[a)] \textbf{Catégories des flèches de}~$\cal{E}$. Soit
$\cal{E}$ une cat\'egorie. D\'esignons par $\mathbf{\Delta^1}$ la
cat\'egorie associ\'ee \`a l'ensemble totalement ordonn\'e
\`a deux \'el\'ements $[0,1]$;
elle a donc deux objets 0 et 1, et en plus des deux morphismes
identiques une fl\`eche $(0,1)$ de source 0 et but 1. Soit
$$
\CatFl(\cal{E})=\SheafHom(\mathbf{\Delta^1},\cal{E})
$$
on l'appelle la \emph{catégorie des flèches de} $\cal{E}$.
L'objet 1 de~$\mathbf{\Delta^1}$ d\'efinit un foncteur canonique,
appel\'e \emph{foncteur-but}
$$
\CatFl(\cal{E})\to \cal{E}
$$
(le foncteur d\'efini par l'objet 0 de~$\mathbf{\Delta^1}$ est
appel\'e \emph{foncteur-source}). Pour tout objet $S$ de~$\cal{E}$,
la cat\'egorie-fibre $\CatFl(\cal{E})_S$ est canoniquement isomorphe
\`a la cat\'egorie $\cal{E}_{/S}$ des objets de~$\cal{E}$
au-dessus de~$S$.

Consid\'erons un morphisme $f\colon T\to S$ dans $\cal{E}$, alors il
lui correspond un foncteur canonique
$$
f_*\colon \cal{E}_{/T}=\cal{F}_T\to \cal{E}_{/S}=\cal{F}_S
$$
et un isomorphisme fonctoriel
$$
\Hom_S(f_*(\eta),\xi)\isomto\Hom_f(\eta,\xi)
$$
qui fait donc de~$f_*$ un foncteur image directe pour $f$ dans
$\cal{F}$. On a d'ailleurs ici
$$
(\id_S)_*=\id_{\cal{F}_S},\quad (fg)_*=f_*g_*,\quad
c^{f,g}=\id_{(fg)},
$$
\ie $\cal{F}$ est muni d'un co-scindage sur~$\cal{E}$. A fortiori,
$\cal{F}$ est co-fibr\'ee sur~$\cal{E}$. Notons maintenant que
l'ensemble des morphismes dans $\cal{F}$ est en correspondance
biunivoque avec l'ensemble des diagrammes carr\'es commutatifs dans
$\cal{E}$.
$$
\xymatrix{{X}\ar[d]_-u&{Y}\ar[l]_-{f'}\ar[d]^-v\\{S}&{T}\ar[l]_-{f}}
$$
Par
d\'efinition, le morphisme en question est cart\'esien si le
carr\'e est cart\'esien dans $\cal{E}$, \ie s'il fait de~$Y$ un
produit fibr\'e de~$X$ et $T$ sur~$S$. Le foncteur image inverse
$f^*$
existe donc si et seulement si pour tout objet $X$ sur~$S$, le produit
fibr\'e $X\times_ST$ existe. Il r\'esulte de~\ref{VI.10.1} que si
le produit de deux objets sur un troisi\`eme existe toujours dans~$\cal{E}$, \ie si $\cal{F}$ est pr\'efibr\'ee sur~$\cal{E}$,
alors $\cal{F}$ est m\^eme fibr\'ee sur~$\cal{E}$.

\medskip
\item[b)] \textbf{Catégorie des préfaisceaux ou faisceaux sur
des espaces variables}

Soit $\cal{E}=\mathbf{Top}$ la cat\'egorie des espaces topologiques.
Si $T$ est un espace topologique, nous noterons $\cal{U}(T)$ la
cat\'egorie des ouverts de~$T$, o\`u les morphismes sont les
applications d'inclusion. Si $\cal{C}$ est une cat\'egorie, un
foncteur $\cal{U}(T)^\circ\to \cal{C}$ s'appelle un
\emph{préfaisceau} sur~$T$ \`a valeurs dans $\cal{C}$, et un
\emph{faisceau} s'il satisfait une condition d'exactitude \`a gauche
que nous ne r\'ep\'etons pas ici. La \emph{catégorie
$\cal{P}(T)$ des préfaisceaux sur~$T$ à valeurs dans
$\cal{C}$}, est par d\'efinition la cat\'egorie
$\SheafHom(\cal{U}(T)^\circ,\cal{C})$, et la cat\'egorie
$\cal{F}(T)$ des faisceaux sur~$T$ \`a valeurs dans $\cal{C}$ est la
sous-cat\'egorie pleine dont les objets sont les objets de
$\SheafHom(\cal{U}(T)^\circ,\cal{C})$ qui sont des faisceaux. Si
$f\colon T\to S$ est un morphisme dans $\cal{E}$, \ie une
application continue d'espaces topologiques, il lui correspond par
l'application croissante $U\mto f^{-1}(U)$ un foncteur
$\cal{U}(S)\to \cal{U}(T)$, d'o\`u un foncteur
$$
f_*\colon\SheafHom(\cal{U}(T)^\circ,\cal{C})\to
\SheafHom(\cal{U}(S)^\circ,\cal{C})
$$
appel\'e \emph{foncteur image directe de préfaisceaux par}
$f$. On voit aussit\^ot que l'image directe d'un faisceau est un
faisceau, donc le foncteur
$f_*\colon \cal{P}(T)\to\cal{P}(S)$
induit un foncteur, \'egalement not\'e
$f_*\colon\cal{F}(T)\to\cal{F}(S)$.
On v\'erifie de plus trivialement (par l'associativit\'e de la
composition des foncteurs) qu'on a, pour une deuxi\`eme application
continue $g\colon U\to T$, l'identit\'e
$$
(gf)_* =g_*f_*,\qquad
\text{de même }\ (\id_S)_*=\id_{\cal{P}(S)}.
$$
De cette fa\c con, on a obtenu un foncteur
$$
S\mto \cal{P}(S)
$$
\resp $$
S\mto \cal{F}(S)
$$
de
$\cal{E}$ dans~$\Cat$. En fait, nous nous int\'eressons au
foncteur correspondant
$$
S\mto \cal{P}(S)^\circ,\qquad\text{\resp~}S\mto
\cal{F}(S)^\circ.
$$
Il d\'efinit une cat\'egorie co-fibr\'ee, et m\^eme
co-scind\'ee, sur la cat\'egorie des espaces topologiques qu'on
appelle la \emph{catégorie co-fibrée des préfaisceaux}
(\resp \emph{faisceaux}) \emph{à valeurs dans} $\cal{C}$
(sous-entendu: sur des espaces variables). Explicitant la construction
du \No \ref{VI.8}, on voit qu'un morphisme d'un pr\'efaisceau $B$ sur~$T$
dans un pr\'efaisceau $A$ sur~$S$ est un couple $(f,u)$ form\'e
d'une application continue de~$T$ dans $S$, et d'un morphisme $u\colon
A\to f_*(B)$ dans la cat\'egorie $\cal{P}(S)$. Cette description
vaut \'egalement pour les morphismes de faisceaux, $\cal{F}$
\'etant une sous-cat\'egorie pleine de~$\cal{P}$.

Dans les cas les plus importants, la cat\'egorie $\cal{P}$ et la
cat\'egorie $\cal{F}$ au-dessus de~$\cal{E}$ sont aussi des
cat\'egories fibr\'ees, \ie pour toute application continue, les
foncteurs image directe $\cal{P}(T)\to \cal{P}(S)$ et
$\cal{F}(T)\to\cal{F}(S)$ ont un foncteur adjoint, qui est alors
not\'e $f^*$ et appel\'e foncteur image inverse de
pr\'efaisceaux \resp foncteur image inverse de faisceaux, par
l'application continue $f$. Ce foncteur existe par exemple si
$\cal{C}=\Ens$. On peut montrer que le foncteur $f^*\colon
\cal{P}(S)\to \cal{P}(T)$ existe chaque fois que dans $\cal{C}$ les
limites inductives (relatives \`a des diagrammes dans l'Univers
consid\'er\'e) existent. La question est moins facile pour
$\cal{F}$; on notera en effet que (m\^eme dans le cas
$\cal{C}=\Ens$) l'image inverse d'un pr\'efaisceau qui est un
faisceau n'est en g\'en\'eral pas un faisceau, en d'autres termes le
foncteur image inverse de faisceau n'est pas isomorphe au foncteur
induit par le foncteur image inverse de pr\'efaisceaux (malgr\'e
la notation commune $f^*$). Ainsi, $\cal{F}$ est une
sous-cat\'egorie co-fibr\'ee de~$\cal{P}$, mais pas une
sous-cat\'egorie fibr\'ee, \ie \emph{le foncteur d'inclusion
$\cal{F}\to\cal{P}$ n'est pas fibrant}.

La cat\'egorie co-fibr\'ee $\cal{P}$ peut se d\'eduire d'une
cat\'egorie co-fibr\'ee (ou plut\^ot fibr\'ee) plus
g\'en\'erale, obtenue ainsi. Pour toute cat\'egorie $\cal{U}$ (dans
l'Univers fix\'e), on pose
$$
\cal{P}(\cal{U})=\SheafHom(\cal{U},\cal{C})
$$
et
on note que $\cal{U}\mto\cal{P}(\cal{U})$ est de fa\c con
naturelle un foncteur contravariant en~$\cal{U}$, de la cat\'egorie
$\Cat$ dans $\Cat$. Il d\'efinit donc une cat\'egorie scind\'ee
au-dessus de~$\cal{E}=\Cat$, que nous noterons $\Cat_{/\!/\cal{C}}$. Les
objets de cette cat\'egorie sont les couples $(\cal{U},p)$ d'une
cat\'egorie $\cal{U}$ et d'un foncteur $p\colon\cal{U}\to\cal{C}$,
et un morphisme de~$(\cal{U},p)$ dans $(\cal{V},q)$ est
essentiellement un couple $(f,u)$, o\`u $f$ est un foncteur
$\cal{U}\to\cal{V}$ et $u$ un homomorphisme de foncteurs $u\colon p\to
qf$. Nous laissons au lecteur le soin d'expliciter la composition des
morphismes dans $\Cat_{/\!/\cal{C}}$. Le foncteur-projection
$$
\cal{F}=\Cat_{/\!/\cal{C}}\to\cal{E}=\Cat
$$
associe au couple $(\cal{U},p)$ l'objet $\cal{U}$; la
cat\'egorie-fibre en $\cal{U}$ est la cat\'egorie
$\SheafHom(\cal{U},\cal{C})$ (\`a isomorphisme pr\`es). Lorsque
dans $\cal{C}$ les limites inductives existent, on montre facilement
que la cat\'egorie fibr\'ee $\Cat_{/\!/\cal{C}}$ sur~$\Cat$ est
\'egalement co-fibr\'ee sur~$\Cat$, \ie on peut d\'efinir la
notion d'\emph{image directe d'un foncteur} $p\colon\cal{U}\to\cal{C}$
par un foncteur $f\colon\cal{U}\to\cal{V}$. La cat\'egorie des
pr\'efaisceaux se d\'eduit de la cat\'egorie fibr\'ee
pr\'ec\'edente par le changement de base
$$
\mathbf{Top}^\circ\to\Cat
$$
(foncteur $S\mto\cal{U}(S)$ d\'efini plus haut), ce qui donne une
cat\'egorie fibr\'ee sur~$\mathbf{Top}^\circ$, et en passant \`a
la cat\'egorie oppos\'ee, on obtient la cat\'egorie
co-fibr\'ee $\cal{P}$ des pr\'efaisceaux au-dessus de
$\mathbf{Top}$. La notion d'image inverse d'un foncteur correspond
\`a celle d'image directe de pr\'efaisceau, la notion d'image
directe d'un foncteur \`a celle d'image inverse d'un
pr\'efaisceau.

\medskip
\item[c)] \textbf{Objets à opérateurs au-dessus d'un objet
à opérateurs}

Soit $\cal{F}$ une cat\'egorie sur~$\cal{E}$, et soit $S$ un objet de
$\cal{E}$ o\`u un groupe $G$ op\`ere, \`a gauche pour fixer les
id\'ees. Cet objet \`a op\'erateurs peut s'interpr\'eter comme
correspondant \`a un foncteur $\lambda\colon\cal{E}'\to \cal{E}$ de
la cat\'egorie (\`a un seul objet, ayant $G$ comme groupe
d'endomorphismes) $\cal{E}'$ d\'efinie par $G$, dans la
cat\'egorie $\cal{E}$, et d\'efinie donc par changement de base
une cat\'egorie $\cal{F}'$ au-dessus de~$\cal{E}'$, qui est
fibr\'ee \resp co-fibr\'ee lorsque $\cal{F}$ l'est
sur~$\cal{E}$. Une
section de~$\cal{E}'$ sur~$\cal{F}'$ (n\'ecessairement
cart\'esienne, car $\cal{E}'$ est un groupo\"ide, et tout
isomorphisme dans $\cal{F}'$ est cart\'esien en vertu
de~\ref{VI.6.12}), peut aussi s'interpr\'eter comme un
$\cal{E}$-foncteur $\cal{E}'\to\cal{F}$ au-dessus de~$\lambda$, ou
aussi comme un objet \`a op\'erateurs $\xi$ dans $\cal{F}$
\og au-dessus\fg de l'objet \`a op\'erateurs $S$.

\medskip
\item[d)] \textbf{Couples de foncteurs adjoints quasi-inverses;
autodualités}

Lorsque la cat\'egorie-base $\cal{E}$ est r\'eduite \`a deux
objets $a,b$ et, en plus des fl\`eches identiques, \`a deux
isomorphismes $f\colon a\to b$ et $g\colon b\to a$ inverse l'un de
l'autre (\ie $\cal{E}$ est un groupo\"ide connexe rigide avec deux
objets), une cat\'egorie cliv\'ee normalis\'ee sur~$\cal{E}$ est
essentiellement la m\^eme chose que le syst\`eme form\'e par
deux cat\'egories $\cal{F}_a$ et $\cal{F}_b$ et un \emph{couple de
foncteurs adjoints} $G\colon \cal{F}_a\to\cal{F}_b$ et $F\colon
\cal{F}_b\to\cal{F}_a$, qui soient des \'equivalences de
cat\'egories (donc quasi-inverses l'un de l'autre). On prendra pour
$\cal{F}_a$ et $\cal{F}_b$ les cat\'egories fibres de~$\cal{F}$,
pour $F$ et $G$ les foncteurs $f^*$ et $g^*$, et les deux
isomorphismes
$$
u\colon FG\isomto\id_{\cal{F}_a}\qquad v\colon
GF\isomto\id_{\cal{F}_b}
$$
sont $c_{g,f}$ et $c_{f,g}$. Les deux conditions usuelles de
compatibilit\'e entre $u$ et $v$ ne sont autres que la condition
\ref{VI.7.4}~$B)$ pour les compos\'es $fgf$ et $gfg$. Il est facile
de montrer que ces conditions suffisent \`a impliquer qu'on a bien
un pseudo-foncteur $\cal{E}^\circ\to\Cat$.

Un cas int\'eressant est celui o\`u l'on a
$$
\cal{F}_b=\cal{F}_a^\circ,\quad G=F^\circ,\quad v=u^\circ.
$$
On appelle \emph{autodualité}
\index{autodualit\'e|hyperpage}dans une cat\'egorie $\cal{C}$, la
donn\'ee d'un foncteur $D\colon \cal{C}\to\cal{C}^\circ$, et d'un
isomorphisme $u\colon DD^\circ\isomto\id_{\cal{C}}$, tels que $u$ et
l'isomorphisme $u^\circ\colon D^\circ D\isomto\id_{\cal{C}^\circ}$
fassent de~$(D,D^\circ)$ un couple de foncteurs adjoints,
(n\'ecessairement quasi-inverses l'un de l'autre). Cette condition
s'\'ecrit:
$$
D(u(x))=u(D(x))\qquad \text{pour tout $x\in\Ob(\cal{C})$.}
$$
\item[e)]
\textbf{Catégories au-dessus d'une catégorie
discrète $\cal{E}$.} On dit que $\cal{E}$ est une
\emph{catégorie discrète} si toute fl\`eche y est une
fl\`eche identique, de sorte que $\cal{E}$ est d\'efini \`a
isomorphisme unique pr\`es par la connaissance de l'ensemble
$I=\Ob(\cal{E})$. La donn\'ee d'une cat\'egorie $\cal{F}$
au-dessus de~$\cal{E}$ \'equivaut donc (\`a isomorphisme unique
pr\`es) \`a la donn\'ee d'une famille de cat\'egories
$\cal{F}_i$ ($i\in I$), les cat\'egories fibres. Toute cat\'egorie
$\cal{F}$ sur~$\cal{E}$ est fibr\'ee, tout $\cal{E}$-foncteur
$\cal{F}\to\cal{G}$ est cart\'esien, on a un isomorphisme canonique
$$ \SheafHom_{\cal{E}/-}(\cal{F},\cal{G}) \isomto
\underset{i}\prod\SheafHom(\cal{F}_i,\cal{G}_i).
$$
En particulier, on obtient
$$
\bf{\Gamma}(\cal{F}/\cal{E})=\varprojLim\cal{F}/\cal{E}\isomto\underset{i}\prod \cal{E}_i.
$$
\item[f)] \textbf{Supposons que $\cal{E}$ ait exactement deux objets
$S$ et $T$,} et en plus des \textbf{morphismes identiques, un
morphisme }$f\colon T\to S$. Alors une cat\'egorie $\cal{F}$
au-dessus de~$\cal{E}$ est d\'efinie, \`a $\cal{E}$-isomorphisme
unique pr\`es, par la donn\'ee de deux cat\'egories $\cal{F}_S$
et $\cal{F}_T$ et d'un bifoncteur $H(\eta,\xi)$ sur~$\cal{F}_T^\circ\times\cal{F}_S$, \`a valeurs dans $\Ens$. En effet,
si~$\cal{F}$ est une cat\'egorie au-dessus de~$\cal{E}$, on lui
associe les deux cat\'egories-fibres $\cal{F}_S$ et~$\cal{F}_T$ et
le bifoncteur $H(\eta,\xi)=\Hom_f(\eta,\xi)$. On laisse au lecteur le
soin d'expliciter la construction en sens inverse. Pour que la
cat\'egorie envisag\'ee soit fibr\'ee (ou pr\'efibr\'ee,
cela revient au m\^eme) il faut et il suffit que le foncteur $H$ soit
repr\'esentable par rapport \`a l'argument $\xi$. Pour qu'elle
soit co-fibr\'ee, il faut et il suffit que $H$ soit
repr\'esentable par rapport \`a l'argument~$\eta$.
\item[g)] Soit $\cal{F}=\cal{C}\times\cal{E}$, consid\'er\'ee
comme cat\'egorie au-dessus de~$\cal{E}$ gr\^ace \`a
$\pr_2$. Alors $\cal{F}$ est fibr\'ee et co-fibr\'ee sur~$\cal{E}$,
et est m\^eme munie d'un scindage et d'un co-scindage canonique,
correspondant au foncteur constant sur~$\cal{E}$, \resp sur~$\cal{E}^\circ$, \`a valeurs dans $\Cat$, de valeur $\cal{C}$. On a
$$
\bf{\Gamma}(\cal{F}/\cal{E})\simeq\SheafHom(\cal{E},\cal{C})
$$
et
$\varprojLim\cal{F}/\cal{E}$ correspond \`a la
sous-cat\'egorie pleine form\'ee des foncteurs
$F\colon\cal{E\to\cal{C}}$ transformant morphismes quelconques en
isomorphismes.
\end{enumerate}

\section{Foncteurs sur une cat\'egorie cliv\'ee}
\label{VI.12}

Soit $\cal{F}$ une cat\'egorie cliv\'ee normalis\'ee
sur~$\cal{E}$. Pour tout objet $S$ de~$\cal{E}$ on d\'esigne par
$$
i_S\colon \cal{F}_S\to\cal{F}
$$
le foncteur d'inclusion. On a donc un homomorphisme fonctoriel, pour
tout morphisme $f\colon T\to S$ dans~$\cal{E}$:
$$
\alpha_f\colon i_T f^*\to i_S,
$$
o\`u $f^*$ est le foncteur changement de base
$\cal{F}_S\to\cal{F}_T$ pour $f$ d\'efini par le clivage. Soit
maintenant
$$
F\colon\cal{F}\to\cal{C}
$$
un foncteur de~$\cal{F}$ dans une cat\'egorie $\cal{C}$, posons,
pour tout~$S\in\Ob(\cal{E})$,
$$
F_S=F\circ i_S\colon\cal{F}_S\to\cal{C}
$$
et pour tout~$f\colon T\to S$ dans~$\cal{E}$,
$$
\varphi_f=F*\alpha_f\colon F_T f^*\to F_S
$$
On a ainsi, \`a tout foncteur $F\colon \cal{F}\to\cal{C}$,
associ\'e une famille $(F_S)$ de foncteurs $\cal{F}_S\to\cal{C}$, et
une famille $(\varphi_f)$ d'homomorphismes de foncteurs $F_T f^*\to
F_S$. Ces familles satisfont aux conditions suivantes:

a) $\varphi_{\id_S}=\id_{F_S}$.

b) Pour deux morphismes $f\colon T\to S$ et $g\colon U\to T$
dans~$\cal{E}$, on a commutativit\'e dans le carr\'e
d'homomorphismes fonctoriels:
$$
\xymatrix@C=2cm{{F_U g^* f^*}\ar[r]^-{F_U*c_{f,g}}\ar[d]_{\varphi_g * f^*}&{F_U(fg)^*}\ar[d]^{\varphi_{fg}}\\{F_T f^*}\ar[r]^-{\varphi_f}&{F_S.}}
$$
La premi\`ere relation est triviale, et la deuxi\`eme relation
s'obtient en appliquant le foncteur~$F$ au diagramme commutatif
$$
\xymatrix@C=2cm{{g^* f^* (\xi)}\ar[r]^-{c_{f,g}(\xi)}\ar[d]_{\alpha_g(f^*(\xi))}&{(fg)^*(\xi)}\ar[d]^{\alpha_{fg}(\xi)}\\{f^*(\xi)}\ar[r]^-{\alpha_f(\xi)}&{\xi}}
$$
pour un objet variable $\xi$ dans~$\cal{F}_S$.

Si $G$ est un deuxi\`eme foncteur $\cal{F}\to\cal{C}$, donnant
naissance \`a des foncteurs $G_S\colon\cal{F}_S\to\cal{C}$ et des
homomorphismes fonctoriels $\psi_f\colon G_Tf\to G_S$, et si $u\colon
F\to G$ est un homomorphisme fonctoriel, alors il lui correspond des
homomorphismes fonctoriels $u*i_S$:
$$
u_S\colon F_S\to G_S
$$
et on constate aussit\^ot que pour tout morphisme $f\colon T\to S$
dans $\cal{E}$, on a commutativit\'e dans les carr\'es
\begin{equation*}
\begin{array}{c}
\xymatrix@C=1.5cm{{F_Tf^*}\ar[r]^-{\phi_f}\ar[d]_{u_T * f^*}&{F_S}\ar[d]^{u_S}\\{G_T f^*}\ar[r]^-{\psi_f}&{G_S}}
\end{array}
\tag*{c)}
\end{equation*}

\begin{proposition}
\label{VI.12.1}
Soit $\cal{H}(\cal{F},\cal{C})$ la cat\'egorie dont les objets sont
les couples de familles $(F_S)$ ($S\in\Ob(\cal{F})$) de foncteurs
$\cal{F}_S\to\cal{C}$, et de familles $(\varphi_f)$
($f\in\Fl(\cal{F})$) d'homomorphismes fonctoriels $F_Tf^*\to F_S$,
satisfaisant les conditions \emph{a) et~b)}, et o\`u les morphismes
sont les familles $(u_S)$ ($S\in\Ob(\cal{F})$) d'homomorphismes
$F_S\to G_S$, v\'erifiant la condition de commutativit\'e
\emph{c)} \'ecrite plus haut, (la composition des morphismes se
faisant par la composition des homomorphismes de foncteurs
$\cal{F}_S\to\cal{C}$). Alors les deux lois explicit\'ees plus haut
d\'efinissent un \emph{isomorphisme} $K$ de la cat\'egorie
$\SheafHom(\cal{F},\cal{C})$ avec la cat\'egorie
$\cal{H}(\cal{F},\cal{C})$.
\end{proposition}

Il est trivial qu'on a bien l\`a un \emph{foncteur} de la
premi\`ere cat\'egorie dans la seconde. Ce foncteur est pleinement
fid\`ele, car pour $F,G$ donn\'es, $\Hom(F,G)\to\Hom(K(F),K(G))$
est trivialement injectif; pour montrer que c'est surjectif, il suffit
de noter que la condition de commutativit\'e c) exprime la
fonctorialit\'e des applications $u(\xi)=u_S(\xi)\colon
F(S)=F_S(\xi)\to G(\xi)=G_S(\xi)$ pour les homomorphismes de la forme
$\alpha_f(\xi)$ dans $\cal{F}$, d'autre part on a la fonctorialit\'e
sur chaque cat\'egorie fibre \ie pour les morphismes dans $\cal{F}$
qui sont des $T$-morphismes ($T\in\Ob(\cal{E})$), d'o\`u la
fonctorialit\'e pour tout morphisme dans $\cal{F}$,
puisqu'un
$f$-morphisme (o\`u $f\colon T\to S$ est un morphisme dans
$\cal{E}$
est de fa\c con unique) un compos\'e d'un morphisme
$\alpha_f(\xi)$ et d'un $T$-morphisme. Il reste donc \`a prouver que
le foncteur~$K$ est bijectif pour les objets. L'argument
pr\'ec\'edent montre d\'ej\`a que $K$ est injectif pour les
objets, reste \`a prouver qu'il est surjectif, \ie que si on part
d'un syst\`eme $(F_S),(\varphi_f)$, satisfaisant a) et b); et si on
d\'efinit une application $\Ob\cal{F}\to\Ob\cal{C}$ par
$$
F(\xi)=F_S(\xi) \quad\text{pour }
\xi\in\Ob\cal{F}_S\subset\Ob\cal{F}
$$
et une application $\Fl(\cal{F})\to\Fl(\cal{C})$ par
$$
F(\alpha_f(\xi)u')=\varphi_f(\xi)\,F_T(u')
$$
pour tout morphisme $f\colon T\to S$ dans $\cal{E}$, tout objet $\xi$
de~$\cal{F}_S$ et tout $T$-morphisme $u'$ de but $f^*(\xi)$, alors on
obtient un \emph{foncteur} $F$ de~$\cal{F}$ dans $\cal{C}$. En effet,
la relation $F(\id_\xi)=\id_{F(\xi)}$ est triviale, il reste \`a
prouver la multiplicativit\'e $F(uv)=F(u)F(v)$ lorsqu'on a un
$f$-morphisme $u\colon\eta\to\xi$ et un $g$-morphisme
$v\colon\zeta\to\nu$, avec $f\colon T\to S$ et $g\colon U\to T$ des
morphismes de~$\cal{E}$. Posant $w=uv$, on aura
$$
u=\alpha_f(\xi)u'\quoi,\quad v=\alpha_g(\eta)v'\quoi,\quad w=\alpha_{fg}(\xi)w'
$$
avec
$$
w'=c_{f,g}(\xi)g^*(u')v'\quad \hbox{(\cf \No \ref{VI.8})}.
$$

Avec ces notations, il faut prouver la commutativit\'e du contour
ext\'erieur du diagramme ci-dessous:
$$
\xymatrix{ F_U(\zeta) \ar@{-<} `u[r] `[rrrrrr]^{F(w')} [rrrrrr]
\ar[rr]^{F_u(v')} \ar[drr]_{F(v)} & & F_Ug^*(\eta)
\ar[rr]^{F_ug^*(u')} \ar[d]^{\varphi_g(\eta)} & & F_Ug^*f^*(\xi)
\ar[rr]^{F_U(c_{f,g}(\xi))}\ar[d]^{\varphi_g(f^*(\xi))}& &
F_U(fg)^*(\xi) \ar[d]^{\varphi_{fg}(\xi)} \\ & & F_T(\eta) \ar@{-<}
`d[r] `[rrrr]_{F(u)} [rrrr] \ar[rr]_{F_T(u')} & & F_Tf^*(\xi)
\ar[rr]_{\varphi_f(\xi)} & & F_S(\xi) }
$$
Or le triangle gauche est commutatif par d\'efinition de~$F(v)$, le
carr\'e m\'edian est commutatif car d\'eduit de l'homomorphisme
$u'\colon \xi\to f(\eta)$ par l'homomorphisme fonctoriel $\alpha_g$,
enfin le carr\'e de droite est commutatif en vertu de la condition
b). La conclusion voulue en r\'esulte.

Supposons maintenant que $\cal{C}$ soit \'egalement une
cat\'egorie cliv\'ee normalis\'ee sur~$\cal{E}$, que nous
appellerons dor\'enavant $\cal{G}$, et que nous nous int\'eressons
aux $\cal{E}$-foncteurs de~$\cal{F}$ dans $\cal{G}$. Si $F$ est un tel
foncteur, il induit des foncteurs
$$
F_S\colon \cal{F}_S\to\cal{G}_S
$$
pour les cat\'egories fibres. D'autre part, pour tout morphisme
$f\colon T\to S$ dans~$\cal{E}$, et tout objet $\xi$ dans~$\cal{F}_S$,
le $f$-morphisme $F(\alpha_f(\xi))$ se factorise de fa\c con unique
par un $T$-morphisme
$$
\varphi_f(\xi)\colon F_T(f^*_{\cal{F}}(\xi))\to
f^*_{\cal{G}}(F_S(\xi))
$$
(o\`u le $\cal{F}$ ou le $\cal{G}$ en indice indique la
cat\'egorie cliv\'ee pour laquelle on prend le foncteur image
inverse), d'o\`u un homomorphisme fonctoriel de foncteurs de
$\cal{F}_S$ dans $\cal{G}_T$:
$$
\varphi_f\colon F_T f^*_{\cal{F}}\to f^*_{\cal{G}}F_S.
$$
Les deux syst\`emes $(F_S)$ et $(\varphi_f)$ satisfont les
conditions suivantes:

a') $\varphi_{\id_S}=\id_{F_S}$.

b') Pour deux morphismes $f\colon T\to S$ et $g\colon U\to T$ dans
$\cal{E}$, on a commutativit\'e dans le diagramme d'homomorphismes
fonctoriels suivant:
$$
\xymatrix{{F_Ug_F^*f_F^*}\ar[rr]^{F_U *c_{f,g}^\cal{F}}\ar[d]_{\varphi_{g^*}f^*_{\cal{F}}}& &{F_U(fg)^*_F}\ar[dd]^{\varphi_{fg}}\\{g^*_{\cal{G}}F_Tf^*_F}\ar[d]_{g^*_{\cal{G}}*\varphi_f}& & \\{g^*_{\cal{G}}f^*_{\cal{G}}F_S}\ar[rr]^{c_{f,g}^\cal{G}*F_S}& &{(fg)^*_{\cal{G}}F_S.}}
$$
Nous en laissons la v\'erification au lecteur, ainsi que
l'\'enonc\'e et la d\'emonstration de l'analogue de la
proposition~\ref{VI.12.1}, impliquant que l'on obtient ainsi une
correspondance biunivoque entre l'ensemble des $\cal{E}$-foncteurs de
$\cal{F}$ dans $\cal{G}$, et l'ensemble des syst\`emes
$(F_S)$,$(\varphi_f)$ satisfaisant les conditions a') et b')
ci-dessus. Bien entendu, dans cette correspondance, les foncteurs
cart\'esiens sont caract\'eris\'es par la propri\'et\'e que
les homomorphismes $\varphi_f$ sont des isomorphismes.

\begin{remarquestar}
Bien entendu, il y a int\'er\^et le plus souvent \`a raisonner
directement sur des cat\'egories fibr\'ees sans utiliser des
clivages explicites, ce qui dispense en particulier de faire appel,
pour la notion simple de~$\cal{E}$-foncteur ou de~$\cal{E}$-foncteur
cart\'esien, \`a une interpr\'etation pesante comme
ci-dessus. C'est pour \'eviter des lourdeurs insupportables, et pour
obtenir des \'enonc\'es plus intrins\`eques,
que nous avons d\^u renoncer \`a partir comme dans \cite{VI.2} de la
notion de cat\'egorie cliv\'ee (appel\'ee \og cat\'egorie
fibr\'ee\fg dans \loccit), qui passe au second rang au profit de
celle de cat\'egorie fibr\'ee. Il est d'ailleurs probable que,
contrairement \`a l'usage encore pr\'epond\'erant maintenant,
li\'e \`a d'anciennes habitudes de pens\'ee, il finira par
s'av\'erer plus commode dans les probl\`emes universels, de ne pas
mettre l'accent sur \emph{une} solution suppos\'ee choisie une fois
pour toutes, mais de mettre toutes les solutions sur un pied
d'\'egalit\'e.
\end{remarquestar}

\begin{thebibliography}{0}{VI.13}
\bibitem[1]{VI.1} A.\, \textsc{Grothendieck}, Sur quelques points d'alg\`ebre
homologique,
T\^ohoku Math.\ J. \textbf{9} (1957) p.\, 119--221.


\bibitem[2]{VI.2} A.\, \textsc{Grothendieck}, Technique de descente et
Th\'eor\`emes d'existence,~I, S\'eminaire Bourbaki~190,
d\'ecembre~1959.
\end{thebibliography}


\refstepcounter{chapter}
\addtocontents{toc}{\medskip\par\noindent\textbf{VII: n'existe pas}}

\chapterspace{-2}
